\chapter{混合模式与辐射通道}\label{chap:hybrid-radiation-channels}

材料的响应谱已经确定可激发自由度，但总电子损失并不等于被收集的光。混合模式、散射及电子整形改变各条通道的权重，需要按具体源与观测窗口分别投影。

\section{混合模式与散射接口}

材料的极化会产生电磁场，电磁场又驱动材料，两者共同形成新的正常振动。如果某个频率范围主要由少数共振贡献，可以通过这些模式的频率、成分和线宽理解电子谱峰。本节先求光与材料的混合模式，再将已建立的散射和响应公式用于光子--电子碰撞以及整形电子对靶的驱动。

\subsection{混合与界面模}

前面的介质响应把材料看成一个频率依赖的 Maxwell 算符，但很多实验还需要回答更具体的问题：当一个电磁共振与一个材料共振频率接近时，系统真正的本征振动还会不会分别是“纯光”和“纯物质”？先固定三个术语。把光--物质耦合人为设为零时各自独立的本征模称为\emph{裸模}；恢复耦合以后重新对角化完整二次哈密顿量得到的本征模称为\emph{混合模}。若调节两个裸频率使它们经过彼此，本来会相交的两条裸频率因为非零耦合而在共振处保持有限最小间距，这种谱线排斥称为\emph{反交叉}。因此只要两个自由度之间存在非零线性耦合，裸模就不再分别守恒，真正的正常模必须从耦合后的二次哈密顿量重新对角化。本征矢同时含有光与物质成分时，这类混合本征模通常称为\emph{极化激元}。

界面束缚模看起来是另一类问题，但数学结构相同。有限维的两个振子通过矩阵行列式寻找极点；连续电磁场则通过满足边界条件的 Maxwell Green 张量寻找极点。先从两个耦合振子求出本征频率与本征矢，就能直接看到反交叉怎样产生；随后再把“求极点”的同一方法用于平面 Maxwell 边值问题：此时极点由 Fresnel 分母或等价的推迟 Green 张量给出。有限维矩阵与连续界面虽然计算形式不同，决定谱峰的对象都是预解算符的极点。

\begin{exbox}[Hopfield 二模模型]
\chapref{chap:general-interaction}的一般二次环境允许在谱上孤立出少数材料自由度。这里取一个类光子模和一个类物质模，作为该联合动力学结构的两模参数化；完整二次模型先保留反旋转项与自极化项，RWA 只在相应无量纲控制量足够小时再施加。

先取一个类光子模 $\hat a_c$ 与一个类物质模 $\hat a_m$，裸频率分别为 $\omega_c,\omega_m>0$，线性耦合为 $g_{\rm cm}$。在不预先丢弃反旋转项时，规范一致的二次模型可写成
\begin{equation}
{
\frac{\hat H_{\rm Hop}}{\hbar}
=\omega_c \hat a_c^\dagger\hat a_c+\omega_m \hat a_m^\dagger\hat a_m
+|g_{\rm cm}|(\hat a_c+\hat a_c^\dagger)(\hat a_m+\hat a_m^\dagger)
+D(\hat a_c+\hat a_c^\dagger)^2 .}
\label{eq:hopfield-full-hamiltonian-dp68}
\end{equation}
这里已通过物质模的相位重定义令 $g_{\rm cm}$ 为非负实数。$D$ 是与光子正则坐标平方相乘的抗磁/自极化频率系数，$[D]={\rm s^{-1}}$；它的数值必须由所采用的微观规范模型确定，不能在强耦合区独立删去。这样的光--物质二次混合模型通常称为 Hopfield 模型。

求完整二次模型的正常模频率及稳定条件，再在旋波近似下求两支混合模式与共振劈裂。

\textbf{解。} 定义
\begin{equation*}
\Omega_c^2\equiv\omega_c(\omega_c+4D).
\end{equation*}
把二次哈密顿量写成两个耦合正则坐标的运动方程，正频正常模满足
\begin{equation}
{
\omega_\pm^2
=\frac{\Omega_c^2+\omega_m^2}{2}
\pm\frac12\sqrt{(\Omega_c^2-\omega_m^2)^2
+16|g_{\rm cm}|^2\omega_c\omega_m}.}
\label{eq:hopfield_full_quadratic_frequencies}
\end{equation}
这两个频率是完整二次模型的极点。稳定性要求二次型正定，特别要求两个 $\omega_\pm^2$ 均为正；量子正常模算符则需用保持玻色辛对易结构的线性正则变换构造。

完整二次光--物质哈密顿量同时包含共振与反旋转项。数守恒模型只有在反旋转支路的能量失配远大于耦合时才成立，其控制量可取
\begin{equation*}
\epsilon_{\rm cr}\equiv\frac{|g_{\rm cm}|}{\omega_c+\omega_m}\ll1,
\qquad
\epsilon_{\rm sp}\equiv\frac{|D|}{\omega_c}\ll1.
\end{equation*}
但旋波近似不能把整个自极化项一起删除，因为
\begin{equation*}
D(\hat a_c+\hat a_c^\dagger)^2
=D\hat a_c^2+D\hat a_c^{\dagger2}
+2D\hat a_c^\dagger\hat a_c+D.
\end{equation*}
其中 $\hat a_c^2$ 与 $\hat a_c^{\dagger2}$ 是快速反旋转项，$2D\hat a_c^\dagger\hat a_c$ 却保持激发数，并在一阶就重整化类光子频率。去掉与动力学无关的常数 $D$，定义
\begin{equation}
{\bar\omega_c\equiv\omega_c+2D,}
\label{eq:hopfield-rwa-renormalized-photon-dp75}
\end{equation}
则数守恒 RWA 哈密顿量为
\begin{equation}
{
\frac{\hat H_{\rm RWA}}{\hbar}
=\bar\omega_c \hat a_c^\dagger\hat a_c+\omega_m \hat a_m^\dagger\hat a_m
+g_{\rm cm}\hat a_c^\dagger\hat a_m+g_{\rm cm}^*\hat a_m^\dagger\hat a_c.}
\label{eq:hopfield-rwa-hamiltonian-dp68}
\end{equation}
在单激发子空间中，对应的 $2\times2$ 厄米矩阵直接给出
\begin{equation}
{
\omega_\pm^{({\rm RWA})}
=\frac{\bar\omega_c+\omega_m}{2}
\pm\sqrt{\frac{(\bar\omega_c-\omega_m)^2}{4}+|g_{\rm cm}|^2}.}
\label{eq:hopfield_rwa_polariton_frequencies}
\end{equation}
在重整化共振条件 $\bar\omega_c=\omega_m$ 处，RWA 劈裂为 $2|g_{\rm cm}|$；远离共振时两支分别恢复重整化的类光子模与裸物质模。混合角满足
$\tan2\theta=2|g_{\rm cm}|/(\bar\omega_c-\omega_m)$。只有在额外近似 $|D|/\omega_c\to0$ 下，才可进一步把 $\bar\omega_c$ 换回 $\omega_c$。下图展示的是保留这一数守恒频移后的 RWA 反交叉；若 $\epsilon_{\rm cr}$ 或 $\epsilon_{\rm sp}$ 不小，应回到\eqnrefs{eq:hopfield_full_quadratic_frequencies}。
\begin{center}
\includegraphics[width=.76\linewidth]{hopfield_avoided_crossing.pdf}
\captionof{figure}{保留数守恒频移的 RWA 反交叉。未耦合类光子频率为 $\bar\omega_c=\omega_c+2D$，共振处最小劈裂为 $2|g_{\rm cm}|$；反旋转修正显著时须使用完整二次模型。}
\label{fig:hopfield-avoided-crossing}
\end{center}

\medskip

取二次哈密顿量，并通过相位重定义令 $g\equiv|g_{\rm cm}|$ 为实数。定义正则正交分量
\begin{equation*}
 \hat q_j=\frac{\hat a_j+\hat a_j^\dagger}{\sqrt2},
 \qquad
 \hat p_j=\frac{\hat a_j-\hat a_j^\dagger}{\ii\sqrt2},
 \qquad [\hat q_j,\hat p_k]=\ii\delta_{jk}.
\end{equation*}
利用
\begin{equation*}
 \hat a_j^\dagger\hat a_j=\frac12(\hat q_j^2+\hat p_j^2-1),
 \qquad
 \hat a_j+\hat a_j^\dagger=\sqrt2\,\hat q_j,
\end{equation*}
忽略与动力学无关的常数后，哈密顿量化为
\begin{equation*}
 \frac{\hat H}{\hbar}
 =\frac{\omega_c}{2}(\hat p_c^2+\hat q_c^2)
 +\frac{\omega_m}{2}(\hat p_m^2+\hat q_m^2)
 +2g\hat q_c\hat q_m+2D\hat q_c^2.
\end{equation*}
Heisenberg 方程逐项给
\begin{align*}
 \dot{\hat q}_c&=\omega_c \hat p_c,&
 \dot{\hat p}_c&=-(\omega_c+4D)\hat q_c-2g\hat q_m,\\
 \dot{\hat q}_m&=\omega_m \hat p_m,&
 \dot{\hat p}_m&=-\omega_m \hat q_m-2g\hat q_c.
\end{align*}
再求一次时间导数并消去 $\hat p_c,\hat p_m$：
\begin{align*}
 \ddot{\hat q}_c
 &=-\omega_c(\omega_c+4D)\hat q_c-2g\omega_c \hat q_m,\\
 \ddot{\hat q}_m
 &=-\omega_m^2\hat q_m-2g\omega_m \hat q_c.
\end{align*}
定义
\begin{equation*}
 \Omega_c^2\equiv\omega_c(\omega_c+4D)
\end{equation*}
并试取正常模拟设 $\hat q_j(t)=Q_j \ee^{-\ii\omega t}$。得到线性方程
\begin{equation*}
 \begin{pmatrix}
 \Omega_c^2-\omega^2 & 2g\omega_c\\
 2g\omega_m & \omega_m^2-\omega^2
 \end{pmatrix}
 \binom{Q_c}{Q_m}=0.
\end{equation*}
非平凡解要求行列式为零：
\begin{align*}
 0
 &=(\Omega_c^2-\omega^2)(\omega_m^2-\omega^2)
 -4g^2\omega_c\omega_m.
\end{align*}
令 $x=\omega^2$，得到普通二次方程
\begin{equation*}
 x^2-(\Omega_c^2+\omega_m^2)x
 +\Omega_c^2\omega_m^2-4g^2\omega_c\omega_m=0.
\end{equation*}
求根并整理判别式：
\begin{align*}
 x_\pm
 &=\frac{\Omega_c^2+\omega_m^2}{2}
 \pm\frac12\sqrt{(\Omega_c^2+\omega_m^2)^2
 -4\Omega_c^2\omega_m^2+16g^2\omega_c\omega_m},\\
 &=\frac{\Omega_c^2+\omega_m^2}{2}
 \pm\frac12\sqrt{(\Omega_c^2-\omega_m^2)^2
 +16g^2\omega_c\omega_m}.
\end{align*}
于是 $x_\pm=\omega_\pm^2$，恢复式。稳定性至少要求两个 $x_\pm>0$；等价地，除迹为正外还需
\begin{equation*}
 \Omega_c^2\omega_m^2-4g^2\omega_c\omega_m>0.
\end{equation*}
抗磁/自极化项不仅改变正常模频率，也参与保证二次哈密顿量的稳定性。真正的量子正常模算符还需由保持玻色辛对易关系的 Bogoliubov 变换构造；普通酉 $2\times2$ 旋转只适用于 RWA 的数守恒子空间。
\end{exbox}

界面问题的逻辑相同，但极点现在来自 Maxwell 边值问题而不是有限维矩阵。考虑 $z=0$ 的两个局域、各向同性、非磁介质 $\varepsilon_1(\omega),\varepsilon_2(\omega)$。对沿界面传播、法向衰减的 TM 场，定义
\begin{equation*}
\alpha_j^2=\kappa_\parallel^2-\varepsilon_j\frac{\omega^2}{c^2},
\qquad
\operatorname{Re}\alpha_j>0.
\end{equation*}
切向 $E$ 与法向 $D$ 的连续性要求
\begin{equation}
{
\frac{\varepsilon_1}{\alpha_1}+\frac{\varepsilon_2}{\alpha_2}=0.}
\label{eq:tm_surface_mode_condition}
\end{equation}
这是应直接求解的表面模方程。无损真空--介质界面中它给出 $\varepsilon_2<-1$ 的束缚 TM 支；准静态极限退化为 $\varepsilon_1+\varepsilon_2=0$。有损体系中 $\kappa_\parallel$ 一般为复数，物理解由推迟/出射支决定。

表面极点与电子耦合同时包含在平面内 Weyl 表示中。对源点与观测点都在介质 1 的情况
\begin{equation}
\mathbf G_{\rm sc}^R(\vb{r},\vb{r}';\omega)
=\int\frac{\dd^2 \kappa_\parallel}{(2\pi)^2}
 \ee^{-\ii\vb{\kappa}_\parallel\cdot(\vb{R}-\vb{R}')}
\mathbf G_{{\rm sc},W}^R(\vb{\kappa}_\parallel;z,z';\omega).
\label{eq:planar_green_weyl_definition}
\end{equation}
令 $\kappa_j^{(0)}=\sqrt{\varepsilon_j}\,\omega/c$、$\kappa_{zj}=\sqrt{(\kappa_j^{(0)})^2-\kappa_\parallel^2}$，并选 $\operatorname{Im}\kappa_{zj}\ge0$ 的推迟支。TE/TM 分解后，散射并矢为
\begin{equation}
{
\mathbf G_{{\rm sc},W}^R
(\vb{\kappa}_\parallel;z,z';\omega)
=\frac{\ii}{2\kappa_{z1}}\ee^{\ii\kappa_{z1}(z+z')}
\left[r_s\,\hat{\vb{s}}\otimes\hat{\vb{s}}
+r_p\,\hat{\vb{p}}_1^+\otimes\hat{\vb{p}}_1^-\right].}
\label{eq:planar_scattering_green_dyadic}
\end{equation}
其中 TM 反射振幅最稳妥地写成复法向波数形式
\begin{equation}
r_p(\kappa_\parallel,\omega)
=\frac{\varepsilon_2\kappa_{z1}-\varepsilon_1\kappa_{z2}}
{\varepsilon_2\kappa_{z1}+\varepsilon_1\kappa_{z2}}
\equiv\frac{N_p}{D_p}.
\label{eq:rp_complex_kz_form}
\end{equation}
因此\eqnrefs{eq:tm_surface_mode_condition} 正是 $D_p=0$ 的极点条件。若 $\kappa_{\rm sp}(\omega)$ 是 $D_p$ 的简单零点，即
\begin{equation*}
 D_p(\kappa_{\rm sp},\omega)=0,
 \qquad
 \partial_{\kappa_\parallel}D_p(\kappa_{\rm sp},\omega)\neq0,
\end{equation*}
则 $r_p=N_p/D_p$ 在该点具有 Laurent 分解
\begin{equation*}
 r_p(\kappa_\parallel,\omega)
 =\frac{\mathcal R_p(\omega)}
 {\kappa_\parallel-\kappa_{\rm sp}(\omega)}
 +r_{p,{\rm reg}}(\kappa_\parallel,\omega),
\end{equation*}
其中
\begin{equation*}
 \mathcal R_p(\omega)
 =\frac{N_p(\kappa_{\rm sp},\omega)}
 {\partial_{\kappa_\parallel}D_p(\kappa_{\rm sp},\omega)},
\end{equation*}
而 $r_{p,{\rm reg}}$ 在极点附近解析。传播长度由 $\operatorname{Im}\kappa_{\rm sp}$ 控制，色散由 $\operatorname{Re}\kappa_{\rm sp}$ 控制，谱权重则由留数 $\mathcal R_p$ 控制；三者是同一个复极点的不同信息。

将 Green 张量投影到电子轨迹后，源的平移动力学直接固定可访问的平行动量。电子以恒速 $v$ 沿 $x$ 方向、在高度 $b$ 处平行界面运动，$\vb{r}_e(t)=(vt,0,b)$。有限时间轨迹先给 $\sinc^2[(\omega-\kappa_xv)T/2]$；只有在 $T\to\infty$ 后才收缩成 $2\pi T\delta(\omega-\kappa_xv)$。因此同步条件在有限作用长度下不是一条无限窄直线，而具有由飞行时间本身决定的 Fourier 宽度。令 $\Delta\equiv\omega-\kappa_xv$。再定义 $x_{1/2}$ 为区间 $(0,\pi)$ 内方程
\begin{equation*}
 \sinc^2x_{1/2}=\frac12
\end{equation*}
的唯一根；数值上 $x_{1/2}=1.391557\ldots$。于是归一化线形及其两个常用宽度严格写成
\begin{equation}
{
\mathcal L_T(\Delta)=\sinc^2\!\left(\frac{\Delta T}{2}\right),
\qquad
\Delta\omega_{\rm zero}=\frac{2\pi}{T},
\qquad
\Delta\omega_{\rm FWHM}=\frac{4x_{1/2}}{T}.}
\label{eq:finite-trajectory-sync-linewidth-dp83}
\end{equation}
这里 $\Delta\omega_{\rm zero}$ 是中心到第一零点的角频率间隔，$\Delta\omega_{\rm FWHM}$ 是主峰全高半宽，数值上 $4x_{1/2}=5.56623\ldots$；用 $L=vT$ 可写成 $\Delta\kappa_{x,{\rm FWHM}}=4x_{1/2}/L$。这给出一个直接实验判据：若两个材料或几何共振在电子同步投影后的角频率间距小于这一 Fourier 宽度，仅靠增加本征模式品质因子并不能使它们在有限轨迹测量中被分辨。长轨迹极限再把该线形收缩为 Dirac $\delta$。因此单位长度的散射电子能量损失数谱为
\begin{equation}
{
\frac1L\frac{\dd P_{\rm sc}}{\dd\omega}
=\frac{\mu_0e^2}{2\pi^2\hbar}
\int_{-\infty}^{\infty}\dd\kappa_y\,
\hat{\vb{x}}\cdot\operatorname{Im}\mathbf G_{{\rm sc},W}^R
\left(\kappa_x=\frac{\omega}{v},\kappa_y;b,b;\omega\right)
\cdot\hat{\vb{x}}.}
\label{eq:planar_eels_mixed_green_per_length}
\end{equation}
该式把源运动学与材料响应分解为两个独立因子：$\kappa_x=\omega/v$ 来自长直轨迹，$\kappa_y$ 仍需积分，而所有结构信息都在 $\mathbf G^R$ 中。

在真空侧的倏逝扇区，$\kappa_{z1}=\ii\alpha_1$，TM 被积函数含有
\begin{equation}
{r_p(\kappa_\parallel,\omega)\ee^{-2\alpha_1b}.}
\label{eq:surface_pole_distance_kernel}
\end{equation}
于是表面谱峰的色散来自 $D_p=0$，权重来自极点留数，碰撞参数依赖来自 $\ee^{-2\alpha_1b}$。三者分别由同一个复极点的位置、留数与真空侧倏逝因子确定。

理想无限长、亚光速的平行轨迹还有一个重要的光锥结论：
\begin{equation}
|\kappa_x|=\frac{\omega}{v}>\frac{\omega}{c},
\qquad \kappa_\parallel\ge |\kappa_x|>\frac{\omega}{c}.
\label{eq:parallel_electron_outside_light_cone}
\end{equation}
因此均匀平面界面上同步的傅里叶支撑全部位于真空光锥外，可以激发倏逝表面响应，却不能仅凭这条无限平移不变轨迹直接产生真空远场光子。有限轨迹、加速和周期结构都会打破这个限制；特别是光栅提供 $G_m$ 后同步条件变为
\begin{equation}
\kappa_x=\frac{\omega}{v}-G_m,
\label{eq:grating_light_cone_momentum_shift}
\end{equation}
从而可以把源谱搬入光锥。这也是 EELS 与 CL 可以共享同一结构极点、却具有不同强度和选择规则的原因。

% -----------------------------------------------------------------------------
% Absolute SI normalization linking EELS, CL/radiation, absorption, and LDOS
% -----------------------------------------------------------------------------

线性电磁响应一旦由推迟 Green 张量确定，同一介质便不需要为不同实验重新建立场方程。差别只在于怎样激发这个响应，以及最后读取哪个投影。

电子能量损失谱（EELS）统计电子转移给环境的总能量。阴极发光（CL）从这些通道中只选择最终进入远场辐射的部分。局域态密度（LDOS）则把电子源换成理想局域偶极源，衡量给定位置可耦合的电磁态。

三者共享同一个 Green 张量，区别只在源、投影和探测通道。确定 Fourier 约定和 SI 单位以后，各自的谱密度归一化也随之唯一确定。

全书采用
\begin{equation*}
\vb{E}(\vb{r},t)
=\int_{-\infty}^{\infty}\frac{\dd\omega}{2\pi}\,
\widetilde{\vb{E}}(\vb{r},\omega)\ee^{-\ii\omega t}.
\end{equation*}
因此时域电场 $[\vb{E}(t)]={\rm V\,m^{-1}}$ 对应
\begin{equation*}
[\widetilde{\vb{E}}(\omega)]={\rm V\,s\,m^{-1}}.
\end{equation*}
时域电流密度满足 $[\vb{J}(t)]={\rm A\,m^{-2}}={\rm C\,s^{-1}m^{-2}}$。同一个 Fourier 变换对时间积分提供一个秒，因此
\begin{equation*}
[\widetilde{\vb{J}}(\vb{r},\omega)]
={\rm C\,m^{-2}}.
\end{equation*}
推迟 Green 张量满足 Maxwell 逆算符方程，已经在一般相互作用章固定为
\begin{equation*}
[\mathbf G^R]={\rm m^{-1}}.
\end{equation*}
这些单位随后唯一固定能量谱与计数谱的归一化。

先定义远场幅度。设探测器位于最终无损真空区域，观察方向为无量纲单位矢量 $\hat{\vb{n}}=\vb{r}/r$，真空波数 $\kappa_0=\omega/c$ 的单位为 ${\rm m^{-1}}$。推迟 Green 张量的 $r\to\infty$ 渐近形式写成
\begin{equation}
{
\mathbf G^R(\vb{r},\vb{r}';\omega)
\xrightarrow[r\to\infty]{}
\frac{\ee^{\ii\kappa_0 r}}{r}
\vb{F}_\infty
(\hat{\vb{n}},\vb{r}';\omega).}
\label{eq:far_field_green_amplitude_definition}
\end{equation}
因为左侧与 $1/r$ 都有 ${\rm m^{-1}}$，远场方向核 $\vb{F}_\infty$ 无量纲。它包含源点 $\vb{r}'$ 到观察方向 $\hat{\vb{n}}$ 的偏振投影、边界散射和材料响应。

对给定频域源电流定义
\begin{equation}
{
\vb{A}_\infty(\hat{\vb{n}},\omega)
\equiv
\int\dd^3 r'\,
\vb{F}_\infty
(\hat{\vb{n}},\vb{r}';\omega)
\cdot\widetilde{\vb{J}}(\vb{r}',\omega).}
\label{eq:far_field_source_amplitude}
\end{equation}
由于 $[\widetilde{\vb{J}}(\omega)]={\rm C\,m^{-2}}$，体积积分给
\begin{equation*}
[\vb{A}_\infty]={\rm C\,m}.
\end{equation*}
由源--Green 关系，远场频域电场为
\begin{equation*}
\widetilde{\vb{E}}_{\rm far}(\vb{r},\omega)
=\ii\mu_0\omega\frac{\ee^{\ii\kappa_0 r}}{r}
\vb{A}_\infty(\hat{\vb{n}},\omega),
\end{equation*}
其单位确实是 ${\rm V\,s\,m^{-1}}$。

真空远场的 Fourier 分量满足 $\widetilde{\vb{H}}_{\rm far}=(\mu_0c)^{-1}\hat{\vb{n}}\times\widetilde{\vb{E}}_{\rm far}$。对正频率用 Poynting 通量积分时间，得到每单位角频率、每单位立体角的辐射\emph{能量谱}
\begin{equation}
{
\frac{\dd^2 W_{\rm rad}}{\dd\omega\,\dd\Omega}
=
\frac{\mu_0\omega^2}{\pi c}
\left|\vb{A}_\infty(\hat{\vb{n}},\omega)\right|^2.}
\label{eq:far_field_radiated_energy_green_si}
\end{equation}
因为 $\dd\omega$ 的单位是 ${\rm s^{-1}}$，该谱的单位为
\begin{equation*}
\left[\frac{\dd^2 W_{\rm rad}}{\dd\omega\dd\Omega}\right]
={\rm J\,s}.
\end{equation*}

若希望把辐射能量换算成“这一频率区间平均包含多少个光子”，定义平均光子数谱
\begin{equation}
{
\frac{\dd^2 \bar N_{\rm rad}}{\dd\omega\,\dd\Omega}
\equiv
\frac{1}{\hbar\omega}
\frac{\dd^2 W_{\rm rad}}{\dd\omega\,\dd\Omega}
=
\frac{\mu_0\omega}{\pi\hbar c}
\left|\vb{A}_\infty\right|^2.}
\label{eq:cl_photon_probability_far_green}
\end{equation}
其单位是秒。对经典给定电流，它首先是平均辐射光子数密度；在单电子弱发射区，若相应频带内的总发射概率远小于一，则同一式子就是一阶单光子发射概率密度。

自由空间为这一归一化提供直接检查。真空远场 Green 核为
\begin{equation*}
\vb{F}_{\infty,0}
=\frac{1}{4\pi}
(\mathbf I-\hat{\vb{n}}\hat{\vb{n}})
\ee^{-\ii\kappa_0\hat{\vb{n}}\cdot\vb{r}'},
\end{equation*}
因此定义空间 Fourier 电流
\begin{equation*}
\widetilde{\vb{J}}(\vb{\kappa},\omega)
\equiv
\int\dd^3 r\,\widetilde{\vb{J}}(\vb{r},\omega)
\ee^{-\ii\vb{\kappa}\cdot\vb{r}},
\qquad
[\widetilde{\vb{J}}]={\rm C\,m}.
\end{equation*}
\eqnrefs{eq:cl_photon_probability_far_green} 退化为
\begin{equation}
{
\frac{\dd^2 \bar N_{\rm rad}^{(0)}}{\dd\omega\,\dd\Omega}
=
\frac{\mu_0\omega}{16\pi^3\hbar c}
\left|
(\mathbf I-\hat{\vb{n}}\hat{\vb{n}})
\widetilde{\vb{J}}(\kappa_0\hat{\vb{n}},\omega)
\right|^2.}
\label{eq:free_space_current_radiation_probability}
\end{equation}
横向投影子删除沿传播方向的电流分量；只有落在真空光锥 $|\vb{\kappa}|=\kappa_0=\omega/c$ 上的横向 Fourier 电流才能把能量输运到无穷远。

远场公式只统计离开控制体的辐射，不能等同于电子的总能量损失。总损失必须从源对自身诱导场所做的功得到；随后再用材料耗散与远场 Poynting 通量把这份总能量分解成互斥去向。频域场先由同一个 Green 张量写成
\begin{equation}
{
\widetilde{\vb{E}}(\vb{r},\omega)
=\ii\mu_0\omega\int\dd^3 r'\,
\mathbf G^R(\vb{r},\vb{r}';\omega)\cdot
\widetilde{\vb{J}}(\vb{r}',\omega).}
\label{eq:eels_source_field_green_dp68}
\end{equation}
采用简写
\begin{equation*}
\langle \widetilde J|\operatorname{Im}G^R|\widetilde J\rangle
\equiv
\int\dd^3 r\dd^3 r'\,
\widetilde{\vb{J}}^{*}(\vb{r},\omega)\cdot
\operatorname{Im}\mathbf G^R(\vb{r},\vb{r}';\omega)
\cdot\widetilde{\vb{J}}(\vb{r}',\omega),
\end{equation*}
该二次型单位为 ${\rm C^2\,m}$。电子或外部给定电流向电磁环境转移的总能量谱定义为
\begin{equation}
{
\frac{\dd W_{\rm loss}}{\dd\omega}
=
\frac{\mu_0\omega}{\pi}
\langle \widetilde J|\operatorname{Im}G^R|\widetilde J\rangle.}
\label{eq:eels_energy_loss_green_si}
\end{equation}
其单位同样为 ${\rm J\,s}$。相应的能量归一化平均激发数谱可定义为
\begin{equation*}
\frac{\dd\bar N_{\rm loss}}{\dd\omega}
\equiv
\frac{1}{\hbar\omega}
\frac{\dd W_{\rm loss}}{\dd\omega}.
\end{equation*}
在弱单电子单量子损失区，$\dd\bar N_{\rm loss}/\dd\omega$ 才可直接解释为 EELS 一阶损失概率密度。

对局域、非磁、被动介质，$\boldsymbol\varepsilon_r(\vb{r},\omega)$ 是无量纲相对介电张量，其厄米虚部 $\operatorname{Im}\boldsymbol\varepsilon_r\succeq0$ 描述不可逆电耗散。相应极化电流为
\begin{equation}
{
\widetilde{\vb{J}}_{\rm mat}
=-\ii\omega\varepsilon_0(\boldsymbol\varepsilon_r-\mathbf I)
\widetilde{\vb{E}}.}
\label{eq:material_polarization_current_dp68}
\end{equation}
它对局域场所做的不可逆功给出材料吸收能量谱
\begin{equation}
{
\frac{\dd W_{\rm abs}}{\dd\omega}
=
\frac{\omega\varepsilon_0}{\pi}
\int_{V_{\rm mat}}\dd^3 r\,
\widetilde{\vb{E}}^{*}(\vb{r},\omega)\cdot
\operatorname{Im}\boldsymbol\varepsilon_r(\vb{r},\omega)
\cdot\widetilde{\vb{E}}(\vb{r},\omega).}
\label{eq:material_absorbed_energy_spectrum_si}
\end{equation}
该式的单位也是 ${\rm J\,s}$。若介质存在磁损耗，需要再加入由 $\operatorname{Im}\boldsymbol\mu_r$ 产生的磁耗散项。

取一个包围所有耗散材料、外边界位于无损真空的控制体。频域 Poynting 定理给出逐频能量守恒：
\begin{equation}
{
\frac{\dd W_{\rm loss}}{\dd\omega}
=
\frac{\dd W_{\rm abs}}{\dd\omega}
+
\int\dd\Omega\,
\frac{\dd^2 W_{\rm rad}}{\dd\omega\,\dd\Omega}.}
\label{eq:eels_radiative_absorptive_energy_balance}
\end{equation}
因此同一总 Green 问题下，辐射分支比最稳妥地定义为\emph{能量通道比}
\begin{equation}
{
\mathcal B_{\rm rad}(\omega)
\equiv
\frac{\displaystyle
\int\dd\Omega\,
\dd^2 W_{\rm rad}/(\dd\omega\dd\Omega)}
{\displaystyle\dd W_{\rm loss}/\dd\omega},
\qquad
0\le\mathcal B_{\rm rad}\le1.}
\label{eq:radiative_branching_ratio_eels_cl}
\end{equation}
$\mathcal B_{\rm rad}$ 无量纲；$1-\mathcal B_{\rm rad}$ 是电吸收分支（若不存在其它能量通道）。若理论先减去自由空间背景或其它参考系统，所得差分谱包含干涉项，可以局部为负；这种差分量不再是独立能量通道，不能套用 $[0,1]$ 的分支比解释。

LDOS 提供一个与电子源无关的同一 Green 归一化检查。\eqnrefs{eq:projected_electric_ldos_definition} 定义的方向投影电学局域态密度单位为
\begin{equation*}
[\rho_{\vb{e}_d}]={\rm s\,m^{-3}},
\end{equation*}
即“每单位角频率、每单位体积”的态密度。真空重合点 Green 张量的有限虚部为
\begin{equation*}
\operatorname{Im}\mathbf G_0^R(\vb{r},\vb{r};\omega)
=\frac{\omega}{6\pi c}\mathbf I,
\end{equation*}
因此任意单位方向上的真空投影 LDOS 为
\begin{equation*}
\rho_{\vb{e}}^{(0)}(\omega)
=\frac{\omega^2}{3\pi^2c^3}.
\end{equation*}
对偶极矩幅值 $d$（单位 ${\rm C\,m}$），代入偶极衰变母\eqnrefs{eq:dipole_decay_projected_ldos} 得
\begin{equation}
{
\Gamma_0
=\frac{\omega^3|d|^2}
{3\pi\varepsilon_0\hbar c^3}.}
\label{eq:vacuum_dipole_decay_ldos_check}
\end{equation}
$\Gamma_0$ 的单位为 ${\rm s^{-1}}$。因此 EELS 的源做功能量、CL 的远场能流和 LDOS 的偶极衰变率虽然实验读出不同，却共享同一个 Green 张量和同一 SI/Fourier 归一化。

\subsection{非迟滞极限：静电反射与闭合损失谱}
\label{chap:feqo-eels}

完整 Fresnel 积分包含传播与倏逝两类场。为了检验归一化并分离材料极点和轨迹权重，下面在同一界面上取非迟滞极限。沿用界面 $z=0$、电子沿 $x$ 运动的坐标；它与电势计算中常用的法向 $x$、运动方向 $z$ 只是坐标互换。

\begin{proposition}[半空间的静电反射核]
设真空在 $z>0$，材料在 $z<0$，电子轨迹为 $(vt,0,b)$。以 $\ee^{-\ii\omega t}$ 为时间约定，并对界面方向作 Fourier 变换，记 $\kappa=(\kappa_x^2+\kappa_y^2)^{1/2}$。在贡献波数上，要求真空侧 $\omega/(c\kappa)\ll1$、材料侧 $|\sqrt\varepsilon|\omega/(c\kappa)\ll1$；两侧法向波数才可同时约化为静电衰减常数。忽略相应迟滞修正后，诱导电势为
\begin{equation}
\widetilde\varphi_{\rm ind}(z,\boldsymbol\kappa,\omega)
=\frac{q_{\mathrm e}}{2\varepsilon_0v\kappa}
r_{\rm es}(\omega)\ee^{-\kappa(z+b)}
2\pi\delta(\kappa_x-\omega/v),
\qquad r_{\rm es}=\frac{1-\varepsilon}{1+\varepsilon}.
\label{eq:feqo-planar-image-potential}
\end{equation}
这里 $r_{\rm es}$ 是电势反射系数；它与本章 TM 电场振幅在静电极限的符号相反。
\end{proposition}
\begin{proof}
电子电荷密度的时间变换为 $q_{\mathrm e}\delta(y)\delta(z-b)\ee^{\ii\omega x/v}/v$。界面 Fourier 变换后，Poisson 方程成为
\[
(\partial_z^2-\kappa^2)\widetilde\varphi=-\widetilde\rho/\varepsilon_0.
\]
点源导数跳跃给出自由项系数 $1/(2\kappa)$。源和界面之间的自由项正比于 $\ee^{\kappa z}$，反射项正比于 $r\ee^{-\kappa z}$，材料中的衰减项正比于 $t\ee^{\kappa z}$。

电势连续与无自由面电荷分别给
\[
1+r=t,\qquad 1-r=\varepsilon t.
\]
消去 $t$ 得 $r=r_{\rm es}$。当 $\varepsilon=1$ 时反射场消失；被动材料满足
\[
-\operatorname{Im}r_{\rm es}
=\frac{2\operatorname{Im}\varepsilon}{|1+\varepsilon|^2}\ge0.
\]
\end{proof}

\begin{theorem}[单位长度非迟滞损失谱]
在长直轨迹、线性响应和弱单次损失条件下，非迟滞概率密度为
\begin{align}
\frac1L\frac{\dd P}{\dd\omega}
&=\frac{q_{\mathrm e}^2}{2\pi\varepsilon_0\hbar v^2}
\int_{-\infty}^{\infty}\frac{\dd\kappa_y}{2\pi}
\frac{\ee^{-2b\sqrt{\kappa_y^2+(\omega/v)^2}}}
{\sqrt{\kappa_y^2+(\omega/v)^2}}
[-\operatorname{Im}r_{\rm es}],\label{eq:feqo-planar-eels}\\
&=\boxed{\frac{q_{\mathrm e}^2}{2\pi^2\varepsilon_0\hbar v^2}
K_0(2\omega b/v)[-\operatorname{Im}r_{\rm es}].}
\label{eq:feqo-planar-eels-k0}
\end{align}
其单位为 ${\rm s\,m^{-1}}$。近似要求实际贡献的波数满足 $\kappa\gg\max\{1,|\sqrt\varepsilon|\}\omega/c$，不能只用颗粒尺寸或“近场”一词替代该条件。在窄共振处还须检查被忽略修正相对反射分母 $1+\varepsilon$ 是否足够小；完整 Fresnel 积分负责处理该极限不均匀的区域。
\end{theorem}
\begin{proof}
将 $E_{{\rm ind},x}=-\ii\kappa_x\varphi_{\rm ind}$ 代入电子损失功，再使用 $\kappa_x=\omega/v$，得到第一式。它也可由完整 TM 核检验：$\alpha_1\to\kappa$、$r_p\to-r_{\rm es}$，TE 项相对消失，$\mu_0/k_0^2=1/(\varepsilon_0\omega^2)$，恰好恢复相同系数。

令 $\kappa_y=(\omega/v)\sinh u$，则
\[
\int_{-\infty}^{\infty}\frac{\dd\kappa_y}{2\pi}
\frac{\ee^{-2b\kappa}}\kappa
=\frac1\pi\int_0^\infty\ee^{-(2\omega b/v)\cosh u}\dd u
=\frac1\pi K_0(2\omega b/v).
\]
这给出第二式，并由被动性保证概率非负。
\end{proof}

\begin{workedexample}[Drude 峰位、线宽与距离权重]
取 $\varepsilon=1-\omega_p^2/[\omega(\omega+\ii\Gamma_D)]$，则
\[
-\operatorname{Im}r_{\rm es}
=\frac{2\Gamma_D\omega\omega_p^2}
{(2\omega^2-\omega_p^2)^2+4\Gamma_D^2\omega^2}.
\]
写 $\omega_s=\omega_p/\sqrt2$、$\delta=\omega-\omega_s$，在 $|\delta|,\Gamma_D\ll\omega_s$ 下，
\[
-\operatorname{Im}r_{\rm es}
\simeq\frac{\Gamma_D\omega_s}{4\delta^2+\Gamma_D^2}.
\]
材料因子的全半高宽为 $\Gamma_D$；只有 $\Gamma_D|\partial_\omega\log K_0(2\omega b/v)|\ll1$ 时，完整谱峰才近似继承这一线宽。

若 $\hbar\omega_p=9\,\mathrm{eV}$，静电极点为 $6.364\,\mathrm{eV}$。完整迟滞计算在 $v=0.70c$ 时的同步尺度却是
\[
\hbar\omega_{\rm sync}
=9\sqrt{\frac{1-0.70^2}{2-0.70^2}}\,\mathrm{eV}
\simeq5.230\,\mathrm{eV}.
\]
这两个数属于不同近似；下一个完整 Fresnel 例题中的约 $5.249\,\mathrm{eV}$ 谱峰不能由 $1+\varepsilon=0$ 直接解释。

当 $X=2\omega b/v\gg1$ 时，用 $\cosh u=1+u^2/2+O(u^4)$ 作 Laplace 展开，得到
\[
K_0(X)=\sqrt{\frac\pi{2X}}\ee^{-X}[1+O(X^{-1})].
\]
故增加距离 $\Delta b$ 主要压低高频信号 $\ee^{-2\omega\Delta b/v}$。束流有位置分布时，应平均这条概率谱，而不是把平均距离代入指数。
\end{workedexample}

% -----------------------------------------------------------------------------
% Worked example: aloof-electron EELS above a planar Drude half-space
% -----------------------------------------------------------------------------

\begin{exbox}[真空--Drude 半空间]
平面界面散射 Green 张量~\eqnrefs{eq:planar_scattering_green_dyadic} 与单位长度损失通式\eqnrefs{eq:planar_eels_mixed_green_per_length} 已给出一般前向模型。将介质选为 Drude 半空间，求指定直线轨迹下的能量损失谱，并用偏振投影、同步条件和数值积分确定峰位。

取界面为 $z=0$，真空位于 $z>0$。电子在高度 $b>0$ 处沿 $x$ 方向匀速运动，
\begin{equation*}
\vb{r}_e(t)=(vt,0,b),
\qquad
\vb{v}=v\hat{\vb{x}},
\qquad 0<v<c.
\end{equation*}
\textbf{解。} 长直轨迹的时间 Fourier 积分已经在正文中给出同步条件
\begin{equation*}
\kappa_x=\frac{\omega}{v}.
\end{equation*}
仍需积分的横向波数为 $\kappa_y$。记
\begin{equation*}
\kappa_\parallel=\sqrt{\kappa_x^2+\kappa_y^2},
\qquad
\kappa_0=\frac{\omega}{c}.
\end{equation*}
由于 $v<c$，有 $|\kappa_x|=\omega/v>\kappa_0$，故整个同步支撑位于真空光锥外。于是
\begin{equation*}
\kappa_{z1}=\ii\alpha_1,
\qquad
\alpha_1=\sqrt{\kappa_\parallel^2-\kappa_0^2}>0.
\end{equation*}
这里 $\alpha_1$ 是真空侧倏逝衰减常数，单位为 ${\rm m^{-1}}$。

先把正文 Green 并矢投影到电子速度方向。真空侧 TE 单位偏振为
\begin{equation*}
\hat{\vb{s}}
=\frac{\hat{\vb{z}}\times\vb{\kappa}_\parallel}{\kappa_\parallel}
=\frac{-\kappa_y\hat{\vb{x}}+\kappa_x\hat{\vb{y}}}{\kappa_\parallel},
\end{equation*}
所以
\begin{equation*}
|\hat{\vb{x}}\cdot\hat{\vb{s}}|^2
=\frac{\kappa_y^2}{\kappa_\parallel^2}.
\end{equation*}
TM 单位偏振取
\begin{equation*}
\hat{\vb{p}}_1^\pm
=\frac{\kappa_\parallel\hat{\vb{z}}
\mp\kappa_{z1}\hat{\vb{\kappa}}_\parallel}{\kappa_0},
\qquad
\hat{\vb{\kappa}}_\parallel
=\frac{\vb{\kappa}_\parallel}{\kappa_\parallel}.
\end{equation*}
因此
\begin{align*}
(\hat{\vb{x}}\cdot\hat{\vb{p}}_1^+)
(\hat{\vb{p}}_1^-\cdot\hat{\vb{x}})
&=-\frac{\kappa_{z1}^2}{\kappa_0^2}
\frac{\kappa_x^2}{\kappa_\parallel^2}\\
&=\frac{\alpha_1^2}{\kappa_0^2}
\frac{\kappa_x^2}{\kappa_\parallel^2}.
\end{align*}
源点和观察点都在 $z=b$，所以 Green 核的公共传播因子化为
\begin{align*}
\frac{\ii}{2\kappa_{z1}}
\ee^{\ii\kappa_{z1}(z+z')}
\bigg|_{z=z'=b}
&=\frac{\ii}{2\ii\alpha_1}\ee^{-2\alpha_1b}\\
&=\frac{\ee^{-2\alpha_1b}}{2\alpha_1}.
\end{align*}
将 TE/TM 两个投影分别乘以 Fresnel 反射振幅 $r_s$、$r_p$，得到电子真正采样的 $xx$ 分量
\begin{equation}
{
G_{{\rm sc},W;xx}^R
=\ee^{-2\alpha_1 b}
\left[
\frac{\kappa_y^2}{2\alpha_1\kappa_\parallel^2}r_s
+\frac{\alpha_1\kappa_x^2}{2\kappa_0^2\kappa_\parallel^2}r_p
\right].}
\label{eq:planar_aloof_gxx_explicit}
\end{equation}
第一项来自 TE 偏振，第二项来自 TM 偏振；平面表面模的极点只出现在相应 Fresnel 振幅中。

把\eqnrefs{eq:planar_aloof_gxx_explicit} 代入正文单位长度 EELS 母\eqnrefs{eq:planar_eels_mixed_green_per_length}，立刻得到完整 Fresnel 积分
\begin{equation}
{
\frac1L\frac{\dd P_{\rm sc}}{\dd\omega}
=\frac{\mu_0e^2}{2\pi^2\hbar}
\int_{-\infty}^{+\infty}\dd\kappa_y\,
\operatorname{Im}\!\left\{
\ee^{-2\alpha_1 b}
\left[
\frac{\kappa_y^2}{2\alpha_1\kappa_\parallel^2}r_s
+\frac{\alpha_1\kappa_x^2}{2\kappa_0^2\kappa_\parallel^2}r_p
\right]\right\}.}
\label{eq:planar_aloof_eels_explicit_fresnel_integral}
\end{equation}
这里 $L$ 是近似平移不变区中的有效轨迹长度。混合 Weyl Green 核满足
\begin{equation*}
[\mathbf G_{{\rm sc},W}^R]={\rm m},
\end{equation*}
因为三维 Green 张量 $[\mathbf G^R]={\rm m^{-1}}$，而 Weyl 逆变换带有 $\dd^2 \kappa_\parallel$。故 $\dd\kappa_y\,G_{{\rm sc},W;xx}^R$ 无量纲，而
\begin{equation*}
\left[\frac{\mu_0e^2}{\hbar}\right]={\rm s\,m^{-1}},
\end{equation*}
从而
\begin{equation*}
\left[\frac1L\frac{\dd P_{\rm sc}}{\dd\omega}\right]
={\rm s\,m^{-1}}.
\end{equation*}
若改用损失能量 $E_{\rm loss}=\hbar\omega$，则
\begin{equation*}
\frac1L\frac{\dd P}{\dd E_{\rm loss}}
=\frac1\hbar\frac1L\frac{\dd P}{\dd\omega}.
\end{equation*}

现在代入具体材料。真空侧 $\varepsilon_1=1$，材料侧取被动 Drude 介电函数
\begin{equation}
{
\varepsilon_2(\omega)
=1-\frac{\omega_p^2}{\omega(\omega+\ii\Gamma_D)}.}
\label{eq:planar_drude_permittivity_example}
\end{equation}
$\omega_p$ 是等离子体角频率，$\Gamma_D$ 是动量弛豫率。在 Fourier 约定 $\ee^{-\ii\omega t}$ 下，$\Gamma_D>0$ 给出正频率处 $\operatorname{Im}\varepsilon_2>0$，因此材料吸收而不主动放大场。

先求一个解析峰位尺度。令 $\Gamma_D=0$，并考察 $\kappa_y=0$ 的 TM 表面支。真空--Drude 界面的表面模色散为
\begin{equation*}
\frac{\kappa_{\rm sp}}{\kappa_0}
=\sqrt{\frac{\varepsilon_2}{1+\varepsilon_2}}.
\end{equation*}
电子同步条件给
\begin{equation*}
\kappa_{\rm sp}=\frac{\omega}{v}
=\frac{\kappa_0}{\beta_{\rm rel}},
\qquad
\beta_{\rm rel}=\frac vc.
\end{equation*}
两式平方并联立：
\begin{align*}
\frac1{\beta_{\rm rel}^2}
&=\frac{\varepsilon_2}{1+\varepsilon_2},\\
1+\varepsilon_2
&=\beta_{\rm rel}^2\varepsilon_2,\\
\varepsilon_2
&=-\frac1{1-\beta_{\rm rel}^2}.
\end{align*}
再代入无损 Drude 形式 $\varepsilon_2=1-\omega_p^2/\omega^2$：
\begin{align*}
1-\frac{\omega_p^2}{\omega^2}
&=-\frac1{1-\beta_{\rm rel}^2},\\
\frac{\omega_p^2}{\omega^2}
&=\frac{2-\beta_{\rm rel}^2}{1-\beta_{\rm rel}^2}.
\end{align*}
因此正频率交叉尺度为
\begin{equation}
{
\omega_{\rm sync}
=\omega_p
\sqrt{\frac{1-\beta_{\rm rel}^2}{2-\beta_{\rm rel}^2}}.}
\label{eq:drude_surface_pole_electron_sync_frequency}
\end{equation}
这不是有限损耗谱峰的精确位置，而是无损、$\kappa_y=0$ 条件下的解析同步尺度。

取数值参数
\begin{equation*}
\hbar\omega_p=\SI{9.0}{eV},
\qquad
\hbar\Gamma_D=\SI{0.08}{eV},
\qquad
v=0.70c,
\qquad
b=\SI{20}{nm}.
\end{equation*}
由\eqnrefs{eq:drude_surface_pole_electron_sync_frequency} 得
\begin{equation*}
\hbar\omega_{\rm sync}=\SI{5.230}{eV}.
\end{equation*}
有限损耗相对这一尺度可用
\begin{equation*}
\epsilon_D=\frac{\Gamma_D}{\omega_{\rm sync}}
\end{equation*}
衡量；$\epsilon_D\ll1$ 时，无损交叉适合作为峰位的初始估计，但完整峰形仍由复 Fresnel 极点、$\kappa_y$ 连续相空间和 $\ee^{-2\alpha_1b}$ 共同决定。

数值上直接积分\eqnrefs{eq:planar_aloof_eels_explicit_fresnel_integral}。上述参数给出谱最大值约在
\begin{equation*}
E_{\rm peak}\simeq\SI{5.249}{eV},
\end{equation*}
峰值约为
\begin{equation*}
\left.\frac1L\frac{\dd P}{\dd E_{\rm loss}}\right|_{\rm peak}
\simeq2.64\times10^{-4}\,{\rm eV^{-1}nm^{-1}}.
\end{equation*}
在弱单电子损失区，可以定义所关心频窗 $\mathcal W_\omega$ 内的总损失概率
\begin{equation*}
\epsilon_{\rm loss}
=L\int_{\mathcal W_\omega}\dd\omega\,
\frac1L\frac{\dd P_{\rm sc}}{\dd\omega}.
\end{equation*}
只有 $\epsilon_{\rm loss}\ll1$ 时，\eqnrefs{eq:planar_aloof_eels_explicit_fresnel_integral} 才能直接解释成单次损失概率密度；更强耦合下需要进一步处理多次能量交换。

\begin{center}
\includegraphics[width=.82\linewidth]{planar_drude_eels.pdf}
\captionof{figure}{真空侧掠射电子对 Drude 半空间的完整迟滞 EELS 谱。实线由完整 Fresnel 核积分；虚线为无损 $\kappa_y=0$ 同步尺度。参数为 $\hbar\omega_p=9.0\,\mathrm{eV}$、$\hbar\Gamma_D=0.08\,\mathrm{eV}$、$v=0.70c$、$b=20\,\mathrm{nm}$。纵轴单位为 ${\rm eV^{-1}nm^{-1}}$。}
\label{fig:planar-drude-eels}
\end{center}

\begin{minipage}{\linewidth}
\hspace*{2em}这个完整例题也显示不同参数进入前向模型的位置不同：$v$ 改变同步线 $\kappa_x=\omega/v$，$b$ 改变高平行动量分量的倏逝权重，$\Gamma_D$ 改变复极点位置和留数。因此三类参数在原则上可以通过谱形和峰位的不同响应加以区分。
\end{minipage}
\end{exbox}


对单振子极性材料，体声子极化激元与表面声子极化激元只是同一 $\varepsilon(\omega)$ 在不同 Maxwell 边界条件下的极点。无阻尼时
\begin{equation*}
\varepsilon(\omega)=\varepsilon_\infty
\frac{\omega_{\rm LO}^2-\omega^2}{\omega_{\rm TO}^2-\omega^2}.
\end{equation*}
体横向模由 $c^2\kappa^2=\varepsilon(\omega)\omega^2$ 决定，得到两条极化激元支；真空界面的 TM 表面支则由\eqnrefs{eq:tm_surface_mode_condition} 决定，在准静态极限满足 $\varepsilon(\omega_s)=-1$，因此
\begin{equation*}
\omega_s^2
=\frac{\varepsilon_\infty\omega_{\rm LO}^2+\omega_{\rm TO}^2}
{\varepsilon_\infty+1}.
\end{equation*}
这说明体模反交叉、表面极点与电子近场损失最终都可追溯到同一个因果材料响应及其不同的边界/源投影。




\subsection{Compton 前向模型}

电磁散射理论给出的 Klein--Nishina 微分截面，给出指定入射光子和电子的一次碰撞概率。实验光束却具有频率和方向分布，探测器也只接收有限范围的出射光。要预言实际计数，需要先对入射分布平均，再对探测范围积分。本节写出这两个步骤，使截面与实验记录直接对应。

先固定单电子实验室系运动学。初始电子静止，入射光子角频率为 $\omega$，散射光子角频率为 $\omega'$，入射与散射方向的夹角为 $\theta\in[0,\pi]$。定义无量纲反冲参数
\begin{equation*}
x\equiv\frac{\hbar\omega}{m_{\mathrm e} c^2}.
\end{equation*}
$x$ 比较入射单光子能量与电子静能。质量壳守恒给
\begin{equation*}
\frac{\omega'}{\omega}
=\frac{1}{1+x(1-\cos\theta)}.
\end{equation*}
因此 $x\ll1$ 时单次散射引起的相对能移至多为 $O(x)$；若希望在全角域采用 Thomson 弹性近似，可取控制量
\begin{equation*}
\epsilon_{\rm Th}\equiv2x\ll1,
\end{equation*}
因为 $0\le1-\cos\theta\le2$。当 $\epsilon_{\rm Th}$ 不小时，必须保留完整 Compton 反冲。

对初态电子静止、初态自旋和光子偏振无偏且末态自旋与偏振不分辨的情形，微观输入取 Klein--Nishina 截面：
\begin{equation}
\frac{\dd\sigma_{\rm KN}}{\dd\Omega_f}
=\frac{r_e^2}{2}\left(\frac{\omega'}{\omega}\right)^2
\left(\frac{\omega}{\omega'}+\frac{\omega'}{\omega}-\sin^2\theta\right),
\qquad r_e=\frac{e^2}{4\pi\varepsilon_0m_{\mathrm e}c^2}.
\label{eq:feqo-compton-microscopic-input}
\end{equation}
其中 $\Omega_f$ 是散射光子的末态立体角，$r_e$ 是经典电子半径，微分截面单位为 ${\rm m^2\,sr^{-1}}$。式中已经包含二体相空间，以及由能量守恒消去 $\omega'$ 产生的雅可比因子。完整的 QED 两图振幅、自旋与偏振求和及相空间推导，见《高等物理》的量子电动力学部分。若初态电子运动，应先在其静止系使用上述截面，再把运动学与通量变换到实验室系。

对非相干入射光束，定义谱角光子通量
\begin{equation*}
\mathscr F_\gamma(\omega,\Omega_i)
\equiv
\frac{\dd^4 N_\gamma}
{\dd t\,\dd A_\perp\,\dd\omega\,\dd\Omega_i}.
\end{equation*}
这里 $A_\perp$ 是垂直于局域入射方向的面积，$\Omega_i$ 是入射方向立体角。因为 $[\dd t]={\rm s}$ 而 $[\dd\omega]={\rm s^{-1}}$，二者在谱通量密度的单位中相消，所以
\begin{equation*}
[\mathscr F_\gamma]={\rm m^{-2}\,sr^{-1}}.
\end{equation*}
它的操作意义是：
\[
\mathscr F_\gamma\,\dd t\,\dd A_\perp\,\dd\omega\,\dd\Omega_i
\]
给出该时空和模式微元内的平均入射光子数。

把探测器的量子效率、几何接受度和已知选择规则合并为无量纲响应
\begin{equation*}
\mathcal Q_{\gamma}^{\rm det}(\omega',\Omega_f),
\qquad
0\le\mathcal Q_{\gamma}^{\rm det}\le1.
\end{equation*}
对一个局域静止电子，探测到的总事件率为
\begin{equation}
{
\Gamma_{\rm det}
=\int\dd\omega\,\dd\Omega_i\,\dd\Omega_f\,
\mathscr F_\gamma(\omega,\Omega_i)
\mathcal Q_{\gamma}^{\rm det}[\omega'(\omega,\theta),\Omega_f]
\frac{\dd\sigma_{\rm KN}}{\dd\Omega_f}.}
\label{eq:compton_detector_forward_rate}
\end{equation}
$[\Gamma_{\rm det}]={\rm s^{-1}}$。这个式子的事件计数结构可以直接在正文中读出。对一个微小模式与立体角单元，
\begin{align*}
\dd N_{\rm det}
={}&\underbrace{\mathscr F_\gamma\,
\dd t\,\dd A_\perp\,\dd\omega\,\dd\Omega_i}_{\text{穿过靶面的入射光子数}}
\underbrace{\frac{1}{\dd A_\perp}
\frac{\dd\sigma_{\rm KN}}{\dd\Omega_f}\dd\Omega_f}_{\text{单电子散射概率}}
\underbrace{\mathcal Q_\gamma^{\rm det}(\omega',\Omega_f)}_{\text{探测条件概率}}.
\end{align*}
面积元 $\dd A_\perp$ 在前两因子之间精确抵消；再除以 $\dd t$ 并对所有入射和末态模式积分，就得到\eqnrefs{eq:compton_detector_forward_rate}。因此入射带宽和角分布只通过 $\mathscr F_\gamma$ 进入，散射后的能量依赖效率只通过 $\mathcal Q_{\gamma}^{\rm det}$ 进入，而微观碰撞只通过 $\dd\sigma_{\rm KN}/\dd\Omega_f$ 进入。

若薄靶包含电子面密度
\begin{equation*}
n_\ee^{(2)}(\vb{R}),
\qquad
[n_\ee^{(2)}]={\rm m^{-2}},
\end{equation*}
其中 $\vb{R}$ 是靶面坐标，则独立单次散射率按 $n_\ee^{(2)}\dd^2 R$ 线性相加。定义典型总 Compton 截面 $\sigma_{\rm C}$，薄靶条件由无量纲光学厚度
\begin{equation*}
\tau_{\rm C}
\equiv n_{e,{\rm char}}^{(2)}\sigma_{\rm C}
\ll1
\end{equation*}
控制。$\tau_{\rm C}$ 达到一阶量级时，光子在离开靶前有不可忽略的再次散射概率，需要用辐射输运或多次散射模型代替线性单碰撞叠加。

\eqnrefs{eq:compton_detector_forward_rate} 假定入射模式构成非相干统计混合。若单光子初态是保持模式间相位的归一化波包
\begin{equation*}
|1_f\rangle
=\sum_\lambda f_\lambda|1_\lambda\rangle,
\qquad
\sum_\lambda|f_\lambda|^2=1,
\end{equation*}
则 $f_\lambda$ 是无量纲复振幅。若不同入射模式 $\lambda$ 可以到达同一个不可区分末态 $f$，总散射振幅必须先相加：
\begin{equation*}
\mathcal A_f^{\rm wp}
=\sum_\lambda f_\lambda\mathcal M_f^{(\lambda)},
\qquad
P_f=|\mathcal A_f^{\rm wp}|^2.
\end{equation*}
交叉项 $f_\lambda f_{\lambda'}^*\mathcal M_\lambda\mathcal M_{\lambda'}^*$ 保留入射波包相位。只有初态密度矩阵在这些模式上已经对角，或实验读出把模式相干完全平均掉时，才能把各模式按 $|f_\lambda|^2$ 在概率层混合。

Thomson 极限由同一微观截面的 $x\to0$ 展开得到。此时 $\omega'/\omega=1+O(x)$，微分截面退化为弹性 Thomson 角分布，\eqnrefs{eq:compton_detector_forward_rate} 相应成为入射光子通量与 Thomson 截面的卷积。高能区 $x\gtrsim1$ 时，$\omega'$ 对 $\theta$ 的强依赖必须继续传入 $\mathcal Q_{\gamma}^{\rm det}(\omega',\Omega_f)$，因为探测器阈值和效率会随散射能量改变。

\subsection{整形电子驱动}\label{subsec:shaped-electron-drive-dp19}

电子还可以充当驱动一个束缚靶的探针。单次碰撞的跃迁由电子产生的场和靶的矩阵元决定；若电子本身是多个动量分量的相干叠加，通向同一靶末态的振幅还可能干涉。让一列电子依次通过时，则要继续比较不同次碰撞的相位。下面先求单电子作用，再求波包相干与电子列同步对靶动力学的影响。

考虑一个 $\Lambda$ 型三能级靶。能量本征态记为 $|1\rangle,|2\rangle,|3\rangle$，本征能量为 $E_1,E_2,E_3$。定义上能级到两个下能级的跃迁角频率
\begin{equation*}
\omega_{3i}\equiv\frac{E_3-E_i}{\hbar}>0,
\qquad i=1,2,
\qquad [\omega_{3i}]={\rm s^{-1}},
\end{equation*}
以及电偶极矩
\begin{equation*}
\vb{d}_{3i}\equiv\langle3|\hat{\vb{d}}|i\rangle,
\qquad
[\vb{d}_{3i}]={\rm C\,m}.
\end{equation*}

单电子沿 $z$ 方向以中心速度 $v_0$ 掠过靶，Lorentz 因子为 $\gamma_0$，横向碰撞参数为
\begin{equation*}
\vb{b}=b\hat{\vb{e}}_\perp,
\qquad [b]={\rm m}.
\end{equation*}
当束缚体系空间尺度 $a_t$ 满足 $a_t/b\ll1$ 时，电子场在靶区域内可按空间坐标展开，第一非零跃迁项是电偶极耦合 $-\hat{\vb{d}}\cdot\vb{E}_e$。

先采用长度为 $L$ 的纵向周期盒归一化。这里 $L$ 的单位为米。用波数 $\kappa$ 而不是物理动量标记电子纵向平面波：
\begin{equation*}
p_z=\hbar \kappa,
\qquad [\kappa]={\rm m^{-1}}.
\end{equation*}
定义电子在一次跃迁中的纵向波数损失
\begin{equation*}
\Delta\kappa\equiv \kappa_{\rm in}-\kappa_{\rm out},
\qquad [\Delta\kappa]={\rm m^{-1}}.
\end{equation*}
匀速电子在靶附近的实验室系电场可写成
\begin{equation}
{
\vb{E}_{e,\perp}(\vb{b},z)
=-\frac{e\gamma_0}{4\pi\varepsilon_0}
\frac{\vb{b}}{(b^2+\gamma_0^2z^2)^{3/2}},
\qquad
E_{e,z}(\vb{b},z)
=-\frac{e\gamma_0}{4\pi\varepsilon_0}
\frac{z}{(b^2+\gamma_0^2z^2)^{3/2}}.}
\label{eq:uniform_electron_boosted_coulomb_field}
\end{equation}
盒归一化的偶极相互作用矩阵元定义为
\begin{equation}
{
 V_{ij}^{(L)}(\Delta\kappa)
=\frac{e}{2\pi\varepsilon_0L\gamma_0}
\left[
|\Delta\kappa|K_1\!\left(\frac{|\Delta\kappa|b}{\gamma_0}\right)
\hat{\vb{e}}_\perp
+\frac{\ii\Delta\kappa}{\gamma_0}
K_0\!\left(\frac{|\Delta\kappa|b}{\gamma_0}\right)
\hat{\vb{e}}_z
\right]\cdot\vb{d}_{ij}.}
\label{eq:feberi_bessel_coulomb_matrix_element}
\end{equation}
$K_0,K_1$ 是无量纲修正 Bessel 函数。方括号内单位为 ${\rm m^{-1}}$，与偶极矩 ${\rm C\,m}$ 收缩后给库仑；前因子 $e/(\varepsilon_0L)$ 单位为伏特，因此
\begin{equation*}
[V_{ij}^{(L)}]={\rm J}.
\end{equation*}
横向 $K_1$ 项和纵向 $K_0$ 项分别来自匀速相对论电子场的横向、纵向 Fourier 分量。

实际散射理论通常使用连续波数归一化
\begin{equation*}
\langle \kappa|\kappa'\rangle=\delta(\kappa-\kappa').
\end{equation*}
因为 $[\delta(\kappa-\kappa')]={\rm m}$，连续态与盒态关系取
\begin{equation*}
|\kappa_n\rangle_L
=\sqrt{\frac{2\pi}{L}}\,|\kappa_n\rangle.
\end{equation*}
于是连续归一化矩阵核 $V_{ij}^{({\rm cont})}(\kappa',\kappa)$ 的单位为 ${\rm J\,m}$，使盒矩阵元恢复焦耳。对近线性色散 $E(\kappa)-E(\kappa-\Delta\kappa)\simeq\hbar v_0\Delta\kappa$，共振波数转移为
\begin{equation*}
u_{ij}\equiv\frac{\omega_{ij}}{v_0},
\qquad [u_{ij}]={\rm m^{-1}}.
\end{equation*}
定义无量纲单碰撞脉冲面积
\begin{equation*}
G_{ij}
\equiv
\frac{2\pi}{\hbar v_0}
 V_{ij}^{({\rm cont})}(u_{ij})
=\frac{L}{\hbar v_0}
 V_{ij}^{(L)}(u_{ij}).
\end{equation*}
$|G_{ij}|$ 比较单次共振碰撞矩阵元与电子飞越产生的作用量尺度；一阶单碰撞展开要求
\begin{equation*}
\epsilon_{G,ij}\equiv|G_{ij}|\ll1.
\end{equation*}
若还要消去靶的其它中间态，则对应耦合与能隙之比还需独立满足 Feshbach/Schrieffer--Wolff 谱隔离条件。

单次碰撞的空间矩阵元只给出“某个波数转移能有多强”；电子纵向态决定不同波数通道能否相干。设连续波数波函数满足
\begin{equation*}
\int\dd \kappa\,|\psi_e(\kappa)|^2=1.
\end{equation*}
因为 $[\dd \kappa]={\rm m^{-1}}$，有
\begin{equation*}
[\psi_e(\kappa)]={\rm m^{1/2}}.
\end{equation*}
定义无量纲纵向波数空间自相关函数
\begin{equation*}
I_e(u)
\equiv\int\dd \kappa\,
\psi_e(\kappa+u)\psi_e^*(\kappa),
\qquad
[I_e]=1.
\end{equation*}
$u$ 是两个电子波数通道的间隔，单位 ${\rm m^{-1}}$。由 Cauchy--Schwarz 不等式有 $|I_e(u)|\le1$。目标跃迁 $i\to3$ 只读取 $u=u_{3i}=\omega_{3i}/v_0$ 处的相干，所以 $I_e(u_{3i})$ 同时给出该波数间隔的相干可见度和相位。

PINEM 相位调制配合自由漂移可以可控地塑造 $I_e(u)$。设调制角频率为 $\omega_{\rm mod}$，定义同步波数
\begin{equation*}
\kappa_{\rm mod}\equiv\frac{\omega_{\rm mod}}{v_0},
\qquad [\kappa_{\rm mod}]={\rm m^{-1}}.
\end{equation*}
调制强度写为无量纲复数
\begin{equation*}
\beta_{\rm mod}=|\beta_{\rm mod}|\ee^{\ii\varphi_{\rm mod}},
\end{equation*}
其中 $\varphi_{\rm mod}$ 是无量纲光学相位。设电子在调制前、后分别自由漂移 $L_s,L_p$，二者单位均为米。在中心波数附近把自由色散保留到二阶：
\begin{equation*}
E(\kappa)
=E_0+\hbar v_0\kappa
+\frac{\hbar^2\kappa^2}{2m_\parallel}.
\end{equation*}
定义二次漂移相位系数
\begin{equation*}
a_s\equiv\frac{\hbar L_s}{2m_\parallel v_0},
\qquad
a_p\equiv\frac{\hbar L_p}{2m_\parallel v_0},
\qquad
[a_s]=[a_p]={\rm m^2}.
\end{equation*}
于是 $a_{s,p}\kappa^2$ 无量纲。最终纵向波函数为
\begin{equation}
{
\begin{aligned}
\psi_e(\kappa)={}&
\ee^{-\ii(\kappa L_p+a_p\kappa^2)}
\sum_{n\in\ZZ}
J_n(2|\beta_{\rm mod}|)\ee^{\ii n\varphi_{\rm mod}}\\
&\times
\ee^{-\ii[(\kappa-n\kappa_{\rm mod})L_s
+a_s(\kappa-n\kappa_{\rm mod})^2]}
\psi_0(\kappa-n\kappa_{\rm mod}).
\end{aligned}}
\label{eq:shaped_electron_pinem_drift_wavefunction}
\end{equation}
PINEM 先生成间隔 $\kappa_{\rm mod}$ 的离散波数副本，漂移二次相位再改变这些副本之间的相对相位。四个控制量
\[
 L_s,\qquad L_p,\qquad |\beta_{\rm mod}|,\qquad \varphi_{\rm mod}
\]
分别调节自由传播相位、调制强度与参考相位，从而共同控制 $I_e(u)$ 的幅度和相位。

\begin{proposition}[相位调制与二次漂移的幺正复合]
设初始纵向波函数属于 $L^2$，调制为实周期正弦相位，自由漂移采用实二次色散。按先漂移、再调制、再漂移的顺序，复合传播精确为\eqnrefs{eq:shaped_electron_pinem_drift_wavefunction}，并保持 $L^2$ 范数。用该二次模型近似真实相对论传播时，另加高阶色散累计相位误差。
\end{proposition}
\begin{proof}
令初始纵向波数波函数为 $\psi_0(\kappa)$，并采用
\[
\int\dd\kappa\,|\psi_0(\kappa)|^2=1.
\]
单频 PINEM 相位调制的 Fourier 展开给出离散波数平移
\[
\psi(\kappa)
\longmapsto
\sum_{n\in\ZZ}
J_n(2|\beta_{\rm mod}|)
\ee^{\ii n\varphi_{\rm mod}}
\psi(\kappa-n\kappa_{\rm mod}).
\]
这里 $n$ 标记净交换的光学量子数，$\kappa_{\rm mod}$ 是对应的纵向波数步长。

自由漂移对每个波数分量只乘一个色散相位。若中心群速度为 $v_0$，长度 $L$ 对应中心飞行时间
$t=L/v_0$。在中心波数附近把能量展开为
\[
E(\kappa)
=E_0+\hbar v_0\kappa
+\frac{\hbar^2\kappa^2}{2m_\parallel}+R_3(\kappa),
\]
其中 $\kappa$ 已取为相对于中心波数的偏移。忽略由相应余项控制量约束的 $R_3$ 后，传播相位为
\begin{align*}
\exp\!\left[-\frac{\ii E(\kappa)L}{\hbar v_0}\right]
={}&\exp\!\left[-\frac{\ii E_0L}{\hbar v_0}\right]
\exp(-\ii\kappa L)
\exp\!\left[-\ii\frac{\hbar L}{2m_\parallel v_0}\kappa^2\right].
\end{align*}
定义
\[
a(L)\equiv\frac{\hbar L}{2m_\parallel v_0}.
\]
第一因子是所有波数分量共有的整体相位，可以从纯态波函数中约去；后两个因子分别表示平移和二次色散。

先让输入波包漂移长度 $L_s$。忽略整体相位后
\[
\psi_s(\kappa)
=\ee^{-\ii[\kappa L_s+a(L_s)\kappa^2]}\psi_0(\kappa).
\]
再施加 PINEM：
\begin{align*}
\psi_{\rm mod}(\kappa)
=\sum_nJ_n(2|\beta_{\rm mod}|)
\ee^{\ii n\varphi_{\rm mod}}
\psi_s(\kappa-n\kappa_{\rm mod}).
\end{align*}
把 $\psi_s$ 的显式形式代入：
\begin{align*}
\psi_{\rm mod}(\kappa)
=\sum_nJ_n(2|\beta_{\rm mod}|)
\ee^{\ii n\varphi_{\rm mod}}
&\exp\!\left\{-\ii[(\kappa-n\kappa_{\rm mod})L_s
+a(L_s)(\kappa-n\kappa_{\rm mod})^2]\right\}\\
&\times\psi_0(\kappa-n\kappa_{\rm mod}).
\end{align*}
最后再漂移 $L_p$，所以末态每个波数分量还要乘
\[
\ee^{-\ii[\kappa L_p+a(L_p)\kappa^2]}.
\]
记 $a_s=a(L_s)$、$a_p=a(L_p)$，按“先漂移 $L_s$、再 PINEM 平移、最后漂移 $L_p$”的时间顺序相乘，得到
\[
\begin{aligned}
\psi_e(\kappa)={}&
\ee^{-\ii(\kappa L_p+a_p\kappa^2)}
\sum_{n\in\ZZ}
J_n(2|\beta_{\rm mod}|)\ee^{\ii n\varphi_{\rm mod}}\\
&\times
\ee^{-\ii[(\kappa-n\kappa_{\rm mod})L_s
+a_s(\kappa-n\kappa_{\rm mod})^2]}
\psi_0(\kappa-n\kappa_{\rm mod}),
\end{aligned}
\]
即正文\eqnrefs{eq:shaped_electron_pinem_drift_wavefunction}。
\end{proof}


$\beta_{\rm mod}$ 采用经典 PINEM 的同一无量纲耦合定义，下标 ``mod'' 只标记它属于电子整形过程。波数自相关 $I_e(u)$ 与纵向密度调制也不是两个独立对象：采用
\begin{equation*}
 \langle z|\kappa\rangle=\frac{1}{\sqrt{2\pi}}\ee^{\ii \kappa z},
 \qquad
 \int\dd z\,|\psi_e(z)|^2=1,
\end{equation*}
直接由 Fourier 变换得到
\begin{equation}
{
 I_e(u)=\int\dd z\,|\psi_e(z)|^2\ee^{-\ii uz}.}
\label{eq:electron_coherence_density_fourier_dp46}
\end{equation}
这里必须区分有限波包的归一化线密度与平均值为一的相对周期密度。将后者写成
$\mathfrak n_e(z)=\sum_dB_d\ee^{\ii d\kappa_{\rm mod}z}$，
其中 $B_0=1$。若非负慢包络 $\eta_e(z)$ 满足 $\int\eta_e\dd z=1$，则物理线密度应写为
\begin{equation*}
 |\psi_e(z)|^2=\frac{\eta_e(z)\mathfrak n_e(z)}{Z_e},\qquad
 Z_e=\int\eta_e(z)\mathfrak n_e(z)\dd z.
\end{equation*}
定义 $\widetilde\eta_e(u)=\int\eta_e(z)\ee^{-\ii uz}\dd z$，代入前式得到有限包络下的结果
\begin{equation*}
 I_e(u)=\frac1{Z_e}\sum_d B_d\widetilde\eta_e(u-d\kappa_{\rm mod}).
\end{equation*}
因此只有在包络 Fourier 宽度远小于 $\kappa_{\rm mod}$、不同谐波的包络尾可忽略时，才有 $Z_e\simeq1$，并在整数谐波处得到
\begin{equation}
{
 I_e(\Delta\ell\,\kappa_{\rm mod})\simeq B_{\Delta\ell}.}
\label{eq:electron_coherence_equals_bunching_dp46}
\end{equation}
在长度为整数个调制周期的均匀盒归一化中，不同谐波严格正交，上式成为等式；有限短波包则应使用保留 $\widetilde\eta_e$ 的求和。这里的 $\eta_e$ 具有 $\mathrm{m^{-1}}$ 量纲，$\mathfrak n_e$、$Z_e$ 与 $I_e$ 均无量纲。
给定 $\beta_{\rm mod}$ 和漂移距离以后，\eqnrefs{eq:pinem_bunching_exact_series_dp46} 决定密度谐波 $B_{\Delta\ell}$，而\eqnrefs{eq:electron_coherence_equals_bunching_dp46} 把同一谐波识别为目标跃迁可读取的电子相干。如果目标频率满足
\begin{equation*}
 \omega_{3i}=\Delta\ell_i\omega_{\rm mod},
 \qquad \Delta\ell_i\in\ZZ,
\end{equation*}
则 $u_{3i}=\Delta\ell_i \kappa_{\rm mod}$，\eqnrefs{eq:wu2026_effective_electron_rabi_frequency} 可直接写成
\begin{equation}
{
 \Omega_{3i}^{(e)}
 \simeq\frac{2G_{3i}}{T}B_{\Delta\ell_i}.}
\label{eq:electron_rabi_from_bunching_dp46}
\end{equation}
于是
\begin{equation*}
 \beta_{\rm mod}
 \longrightarrow B_{\Delta\ell_i}(L_s,L_p)
 \longrightarrow\Omega_{3i}^{(e)}
\end{equation*}
构成从 PINEM 调制到靶 Rabi 频率的定量前向链，因此靶 Rabi 频率由上游 PINEM 调制、自由传播和靶矩阵元共同决定。

再把单次碰撞放入周期电子列。令相邻电子到达靶的时间间隔为
\begin{equation*}
T>0,
\qquad [T]={\rm s}.
\end{equation*}
第 $i$ 条跃迁的单碰撞无量纲相干冲量为 $G_{3i}I_e(u_{3i})$。若电子列与跃迁相位同步，平均到每个周期的相干驱动哈密顿量为
\begin{equation}
{
H_\ee^{(1)}
=\frac{\hbar}{T}
\left[
G_{31}I_e(u_{31})|3\rangle\langle1|
+G_{32}I_e(u_{32})|3\rangle\langle2|
+\mathrm{H.c.}
\right].}
\label{eq:electron_train_average_hamiltonian_dp68}
\end{equation}
采用 Rabi 频率约定
\begin{equation*}
H_\ee^{(1)}
=\frac{\hbar}{2}
\sum_{i=1,2}
[\Omega_{3i}^{(e)}|3\rangle\langle i|+\mathrm{H.c.}],
\end{equation*}
得到
\begin{equation}
{
\Omega_{3i}^{(e)}
=\frac{2G_{3i}I_e(u_{3i})}{T},
\qquad
[\Omega_{3i}^{(e)}]={\rm s^{-1}}.}
\label{eq:wu2026_effective_electron_rabi_frequency}
\end{equation}
$G_{3i}$ 由电子轨迹与偶极取向决定，$I_e$ 由电子纵向相干决定，$1/T$ 则把每次碰撞的无量纲冲量转换成连续时间驱动率；系数 $2$ 来自两种 Rabi 哈密顿量约定的直接比较。

有限电子列只有有限的频率选择性。对第 $i$ 条跃迁，取最接近 $\omega_{3i}T/(2\pi)$ 的整数 $m_i$，定义角频率失谐和每周期相位失配
\begin{equation*}
\Delta_i^{({\rm tr})}
\equiv\omega_{3i}-\frac{2\pi m_i}{T},
\qquad
[\Delta_i^{({\rm tr})}]={\rm s^{-1}},
\end{equation*}
\begin{equation*}
\delta_i\equiv\Delta_i^{({\rm tr})}T
=\omega_{3i}T-2\pi m_i.
\end{equation*}
$\delta_i$ 无量纲。若连续 $N$ 个电子具有相同单次复振幅 $\mathscr A_1^{\rm cell}$，第 $n$ 次事件只多出相位 $n\delta_i$，因此
\begin{equation}
{
\mathscr A_N^{\rm arr}
=\mathscr A_1^{\rm cell}\sum_{n=0}^{N-1}\ee^{\ii n\delta_i}
\equiv\mathscr A_1^{\rm cell} S_N(\delta_i).}
\label{eq:electron_train_phase_sum_dp68}
\end{equation}
把有限几何级数求和后，得到相干因子的闭式
\begin{equation}
{
S_N(\delta_i)
=\sum_{n=0}^{N-1}\ee^{\ii n\delta_i}
=\ee^{\ii(N-1)\delta_i/2}
\frac{\sin(N\delta_i/2)}{\sin(\delta_i/2)}.}
\label{eq:electron-train-coherence-factor-dp32}
\end{equation}
$S_N$ 无量纲。严格同步时 $|S_N|=N$；第一零点满足 $|\delta_i|=2\pi/N$，对应角频率宽度
\begin{equation*}
|\Delta_i^{({\rm tr})}|
\sim\frac{2\pi}{NT}.
\end{equation*}
因此列长 $NT$ 决定傅里叶限制的同步带宽。

\begin{proposition}[有限列 Fourier 相干因子与独立抖动]
设 $N\ge1$ 为整数，单事件包络可 Fourier 变换，相邻事件间隔 $T$。列频谱等于单事件频谱乘 $S_N(\delta)=\sum_{n=0}^{N-1}\ee^{\ii n\delta}$，其归一化模平方为\eqnrefs{eq:electron-train-coherence-factor-dp32}。独立零均值高斯到达误差给 $\mathbb E|S_N^{\rm jit}|^2=N+\ee^{-\omega^2\sigma_t^2}(|S_N|^2-N)$。$N=1$ 时没有相干零点。
\end{proposition}
\begin{proof}
把一次电子事件对目标跃迁产生的有效驱动包络记为 $f_1(t)$。等间隔 $T$ 的有限 $N$ 电子列为
\[
f_N(t)=\sum_{n=0}^{N-1}f_1(t-nT).
\]
Fourier 变换的平移定理给出严格因子化
\begin{align*}
\widetilde f_N(\omega)
&=\widetilde f_1(\omega)
\sum_{n=0}^{N-1}\ee^{+\ii\omega nT}.
\end{align*}
对最接近的同步谐波写 $\omega T=2\pi m+\delta$，整数相位 $\ee^{\ii2\pi mn}=1$，所以列因子正是正文的
\[
S_N(\delta)=\sum_{n=0}^{N-1}\ee^{\ii n\delta}.
\]
乘以 $1-\ee^{\ii\delta}$ 后，中间项逐项消去：
\[
(1-\ee^{\ii\delta})S_N(\delta)=1-\ee^{\ii N\delta}.
\]
再用
\[
1-\ee^{\ii x}=-2\ii \ee^{\ii x/2}\sin\frac{x}{2}
\]
分别处理分子、分母，得到
\[
S_N(\delta)
=\ee^{\ii(N-1)\delta/2}
\frac{\sin(N\delta/2)}{\sin(\delta/2)}.
\]
连续延拓到 $\delta=0$ 给 $S_N(0)=N$。因此归一化相干增强为
\[
\mathcal C_N(\delta)
\equiv\frac{|S_N(\delta)|^2}{N^2}
=\frac1{N^2}
\frac{\sin^2(N\delta/2)}{\sin^2(\delta/2)},
\]
当 $N\ge2$ 时，离同步点最近的零点严格位于 $|\delta|=2\pi/N$，对应
$|\Delta^{({\rm tr})}|=2\pi/(NT)$。

时间抖动也可以直接放进这个和式，而不只写成定性条件。若第 $n$ 个电子的到达误差 $\tau_n$ 独立服从零均值高斯分布，方差为 $\sigma_t^2$，则
\[
S_N^{\rm jit}
=\sum_{n=0}^{N-1}
 \ee^{\ii n\delta}\ee^{\ii\omega\tau_n}.
\]
对强度取系综平均。对 $n\neq m$，独立高斯特征函数给
\[
\left\langle
 \ee^{\ii\omega(\tau_n-\tau_m)}
\right\rangle
=\ee^{-\omega^2\sigma_t^2},
\]
而 $n=m$ 项恒为一，因此
\[
{
\left\langle|S_N^{\rm jit}|^2\right\rangle
=N+\ee^{-\omega^2\sigma_t^2}
\left(|S_N(\delta)|^2-N\right).}
\]
\end{proof}


若电子到达时间存在均方根抖动 $\sigma_t$，第 $i$ 条跃迁的单次相位噪声尺度是 $\omega_{3i}\sigma_t$；若靶的相干时间为 $T_2$，则有效相干电子数不能超过 $T_2/T$。所以 $N^2$ 概率增强要求同时满足
\begin{equation*}
N|\delta_i|\lesssim1,
\qquad
\omega_{3i}\sigma_t\ll1,
\qquad
NT\lesssim T_2.
\end{equation*}

在两条驱动均近共振时，$H_\ee^{(1)}$ 只耦合下能级空间中的一个亮态。定义总 Rabi 尺度
\begin{equation*}
\Omega_B
\equiv\sqrt{|\Omega_{31}^{(e)}|^2+|\Omega_{32}^{(e)}|^2},
\end{equation*}
则与 $|3\rangle$ 正交解耦的归一化暗态为
\begin{equation*}
|D_e\rangle
=\frac{\Omega_{32}^{(e)}|1\rangle
-\Omega_{31}^{(e)}|2\rangle}{\Omega_B}.
\end{equation*}
它满足 $\langle3|H_\ee^{(1)}|D_e\rangle=0$。真实长期稳态还由两次电子碰撞之间的自由开放动力学决定。令 $\mathcal L_0$ 为靶的 Liouville 生成元，单位 ${\rm s^{-1}}$，$\mathscr E_e$ 为一次电子碰撞诱导的无量纲 CPTP 映射，则一个周期的频闪映射是
\begin{equation*}
\mathscr T_{\rm cyc}
\equiv\mathscr E_0\circ\mathscr E_e,\qquad\mathscr E_0\equiv\ee^{T\mathcal L_0}.
\end{equation*}
稳态密度矩阵满足
\begin{equation*}
\mathscr T_{\rm cyc}[\rho_{\rm ss}]=\rho_{\rm ss}.
\end{equation*}
若在迹零扰动子空间上 $\mathscr T_{\rm cyc}$ 的最大本征值模严格小于一，则所有初态扰动都随周期数衰减，固定点唯一。

有限列长改变 Dirichlet 核本身，而到达时间抖动只压低不同电子事件之间的交叉项；两种展宽机制在公式中可直接分开。

\begin{derivation}[平面界面的混合表象与单位长度损失]
先从固定平行动量 $\boldsymbol\kappa_\parallel=(\kappa_x,\kappa_y)$ 的 $z$ 向标量 Green 核开始。介质 1 中它满足
\begin{equation*}
 (\partial_z^2+\kappa_{z1}^2)\mathscr G_z(z,z')
 =-\delta(z-z').
\end{equation*}
在 $z\neq z'$ 的区域，齐次解由推迟边界条件 $\operatorname{Im}\kappa_{z1}\ge0$ 选成
\begin{equation*}
 \mathscr G_z(z,z')
 =A\,\ee^{\ii\kappa_{z1}|z-z'|}.
\end{equation*}
为了确定 $A$，在 $z'\pm\epsilon$ 间积分微分方程：
\begin{align*}
 \int_{z'-\epsilon}^{z'+\epsilon}\dd z\,
 \partial_z^2\mathscr G_z
 +\kappa_{z1}^2
 \int_{z'-\epsilon}^{z'+\epsilon}\dd z\,\mathscr G_z
 &=-1.
\end{align*}
当 $\epsilon\to0$，第二项按 $O(\epsilon)$ 消失，所以
\begin{equation*}
 \partial_z\mathscr G_z\big|_{z'+0}
 -\partial_z\mathscr G_z\big|_{z'-0}=-1.
\end{equation*}
由
\begin{align*}
 \partial_z \ee^{\ii\kappa|z-z'|}
 &=\ii\kappa\,\operatorname{sgn}(z-z')
 \ee^{\ii\kappa|z-z'|}
\end{align*}
可得跳跃
\begin{equation*}
 2\ii\kappa_{z1}A=-1,
 \qquad
 A=\frac{\ii}{2\kappa_{z1}}.
\end{equation*}
这就是 Weyl 表示中每个平行动量通道的标量传播因子。

界面反射路径从源点 $z'$ 到界面再回到观察点 $z$，因此自由空间相位 $\ee^{\ii\kappa_{z1}|z-z'|}$ 变成
\begin{equation*}
 \ee^{\ii\kappa_{z1}(z+z')}.
\end{equation*}
对 TE、TM 两个物理偏振分别乘 Fresnel 反射系数 $r_s,r_p$，并插入对应的偏振二张量，就得到正文\eqnrefs{eq:planar_scattering_green_dyadic}。电子轨迹投影随后决定同步条件与有限长度线宽。

电子轨迹为
\begin{equation*}
 \mathbf r_e(t)=(vt,0,b),
 \qquad -\frac T2\le t\le\frac T2.
\end{equation*}
Weyl 相位含
\begin{equation*}
 \ee^{\ii\boldsymbol\kappa_\parallel\cdot
 [\mathbf R_e(t)-\mathbf R_e(t')]}
 =\ee^{\ii\kappa_xv(t-t')}.
\end{equation*}
频域 Green 泛函还含 $\ee^{\ii\omega(t'-t)}$，因此双时间积分的全部纵向相位组合成
\begin{equation*}
 \ee^{-\ii\Delta(t-t')},
 \qquad
 \Delta\equiv\omega-\kappa_xv.
\end{equation*}
双积分可先因子化：
\begin{align*}
 I_T(\Delta)
 &\equiv
 \int_{-T/2}^{T/2}\dd t\,
 \ee^{-\ii\Delta t}\\
 &=\left.
 \frac{\ee^{-\ii\Delta t}}{-\ii\Delta}
 \right|_{-T/2}^{T/2}\\
 &=\frac{2\sin(\Delta T/2)}{\Delta}\\
 &=T\,\sinc\!\left(\frac{\Delta T}{2}\right).
\end{align*}
于是
\begin{align*}
 &\int_{-T/2}^{T/2}\dd t
 \int_{-T/2}^{T/2}\dd t'\,
 \ee^{-\ii\Delta(t-t')}\\
 &\qquad=I_T(\Delta)I_T^*(\Delta)\\
 &\qquad=T^2\sinc^2\!\left(\frac{\Delta T}{2}\right).
\end{align*}
因此有限轨迹的同步线形不是 $\delta$ 函数，而是正文\eqnrefs{eq:finite-trajectory-sync-linewidth-dp83} 的 $\sinc^2$ 主峰。

第一零点来自
\begin{equation*}
 \sin\frac{\Delta T}{2}=0,
 \qquad
 \frac{|\Delta|T}{2}=\pi,
\end{equation*}
所以从中心到第一零点的角频率间隔是
\begin{equation*}
 \Delta\omega_{\rm zero}=\frac{2\pi}{T}.
\end{equation*}
半高点令
\begin{equation*}
 \sinc^2x=\frac12,
 \qquad x>0.
\end{equation*}
即求方程
\begin{equation*}
 \sin x=\frac{x}{\sqrt2}.
\end{equation*}
主峰内唯一正根为
\begin{equation*}
 x_{1/2}=1.391557\ldots.
\end{equation*}
左右半高点满足
\begin{equation*}
 \Delta_\pm=\pm\frac{2x_{1/2}}{T},
\end{equation*}
所以全高半宽
\begin{align*}
 \Delta\omega_{\rm FWHM}
 &=\Delta_+-\Delta_-\\
 &=\frac{4x_{1/2}}{T}\\
 &=\frac{5.56623\ldots}{T},
\end{align*}
即正文\eqnrefs{eq:finite-trajectory-sync-linewidth-dp83}。

再严格验证长时间极限。对任意在 $\Delta=0$ 附近连续且可积的测试函数 $f(\Delta)$，考虑
\begin{align*}
 J_T
 &\equiv\int_{-\infty}^{\infty}\dd\Delta\,
 f(\Delta)
 T\sinc^2\!\left(\frac{\Delta T}{2}\right).
\end{align*}
令 $x=\Delta T/2$，则 $\dd\Delta=2\dd x/T$：
\begin{align*}
 J_T
 &=2\int_{-\infty}^{\infty}\dd x\,
 f\!\left(\frac{2x}{T}\right)
 \sinc^2x.
\end{align*}
当 $T\to\infty$，对固定 $x$ 有 $f(2x/T)\to f(0)$，且
\begin{equation*}
 \int_{-\infty}^{\infty}\dd x\,\sinc^2x=\pi.
\end{equation*}
因此
\begin{equation*}
 \lim_{T\to\infty}J_T
 =2\pi f(0),
\end{equation*}
也就是分布恒等式
\begin{equation*}
 T\sinc^2\!\left(\frac{\Delta T}{2}\right)
 \xrightarrow[T\to\infty]{\mathcal D'}
 2\pi\delta(\Delta).
\end{equation*}

现在回到 Weyl 动量积分。长轨迹极限把
\begin{equation*}
 \delta(\omega-v\kappa_x)
\end{equation*}
插入 $\kappa_x$ 积分。由 $g(\kappa_x)=\omega-v\kappa_x$、$|g'|=v$，
\begin{equation*}
 \delta(\omega-v\kappa_x)
 =\frac1v
 \delta\!\left(\kappa_x-\frac\omega v\right).
\end{equation*}
因此
\begin{align*}
 \int_{-\infty}^{\infty}\frac{\dd\kappa_x}{2\pi}
 F(\kappa_x,\kappa_y)
 2\pi\delta(\omega-v\kappa_x)
 &=\frac1v
 F\!\left(\frac\omega v,\kappa_y\right).
\end{align*}
双时间积分提供的总因子是 $T$，而概率除以轨迹长度 $L=vT$；再乘上上式 Jacobian $1/v$ 与电流中的两个速度因子 $v^2$，有
\begin{equation*}
 \frac{T}{L}\times\frac1v\times v^2
 =1.
\end{equation*}
所以单位长度谱中不再残留额外的 $T$ 或 $v$ 前因子，$\kappa_x$ 被固定在同步值 $\omega/v$，只剩 $\kappa_y$ 积分。这正得到正文\eqnrefs{eq:planar_eels_mixed_green_per_length}。
\end{derivation}

\begin{derivation}[匀速电子的 Coulomb 场与跃迁矩阵元]
先取给定冲击参数下的直线匀速电子轨迹，并在这一计算层级忽略环境反作用对轨迹的改变。以靶中心为原点，轨迹写成 $(\mathbf b,z)$；相应实验室系 Liénard--Wiechert 电场为

\begin{align*}
\mathbf E_{e,\perp}(\mathbf b,z)
&=-\frac{e\gamma_0}{4\pi\varepsilon_0}
\frac{\mathbf b}{(b^2+\gamma_0^2z^2)^{3/2}},\\
E_{e,z}(\mathbf b,z)
&=-\frac{e\gamma_0}{4\pi\varepsilon_0}
\frac{z}{(b^2+\gamma_0^2z^2)^{3/2}}.
\end{align*}
电偶极相互作用为 $V(z)=-\mathbf d\cdot\mathbf E_e(z)$。盒平面波矩阵元含 Fourier 因子
\[
\frac1L\int_{-\infty}^{\infty}\dd z\,
\ee^{\ii\Delta\kappa z}V(z).
\]
令 $x=\gamma_0z$，则 $\dd z=\dd x/\gamma_0$，指数变成
\[
\ee^{\ii(\Delta\kappa/\gamma_0)x}.
\]
先以 $\kappa_x$ 表示一维傅里叶波数参数，单位 ${\rm m^{-1}}$。从基本积分
\[
\int_{-\infty}^{\infty}\frac{\ee^{\ii\kappa_x x}}{\sqrt{x^2+b^2}}\dd x
=2K_0(|\kappa_x|b)
\]
出发。对 $b$ 求导：
\begin{align*}
\frac{\partial}{\partial b}
\frac1{\sqrt{x^2+b^2}}
&=-\frac{b}{(x^2+b^2)^{3/2}},\\
\frac{\partial}{\partial b}2K_0(|\kappa_x|b)
&=-2|\kappa_x|K_1(|\kappa_x|b),
\end{align*}
所以
\[
\int\dd x\,
\frac{b\ee^{\ii\kappa_x x}}{(x^2+b^2)^{3/2}}
=2|\kappa_x|K_1(|\kappa_x|b).
\]
纵向积分利用
\[
\frac{x}{(x^2+b^2)^{3/2}}
=-\frac{\partial}{\partial x}
\frac1{\sqrt{x^2+b^2}}.
\]
分部积分，边界项在 $|x|\to\infty$ 消失：
\begin{align*}
\int\dd x\,
\frac{x\ee^{\ii\kappa_x x}}{(x^2+b^2)^{3/2}}
&=-\int\dd x\,\ee^{\ii\kappa_x x}
\frac{\partial}{\partial x}
\frac1{\sqrt{x^2+b^2}}\\
&=\ii\kappa_x\int\dd x\,
\frac{\ee^{\ii\kappa_x x}}{\sqrt{x^2+b^2}}\\
&=2\ii\kappa_xK_0(|\kappa_x|b).
\end{align*}
在积分恒等式中取参数 $|\Delta\kappa|/\gamma_0$，再分别乘回横向和纵向偶极分量以及盒归一化 $1/L$，即得\eqnrefs{eq:feberi_bessel_coulomb_matrix_element}。
\end{derivation}

\begin{derivation}[远场振幅与辐射能谱]
真空远场的 Fourier 分量满足

\begin{equation*}
 \widetilde{\mathbf H}_{\rm far}(\omega)
 =\frac1{\mu_0c}\hat{\mathbf n}\times
 \widetilde{\mathbf E}_{\rm far}(\omega),
 \qquad
 \hat{\mathbf n}\cdot\widetilde{\mathbf E}_{\rm far}=0.
\end{equation*}
时域总辐射能为
\begin{align*}
 \frac{\dd W_{\rm rad}}{\dd\Omega}
 &=r^2\int_{-\infty}^{\infty}\dd t\,
 \hat{\mathbf n}\cdot
 [\mathbf E_{\rm far}(t)\times\mathbf H_{\rm far}(t)]\\
 &=\frac{r^2}{\mu_0c}
 \int_{-\infty}^{\infty}\dd t\,|\mathbf E_{\rm far}(t)|^2.
\end{align*}
采用
\begin{equation*}
 \mathbf E(t)=\int\frac{\dd\omega}{2\pi}\,
 \widetilde{\mathbf E}(\omega)\ee^{-\ii\omega t},
\end{equation*}
Parseval 恒等式给
\begin{equation*}
 \int\dd t\,|\mathbf E(t)|^2
 =\int\frac{\dd\omega}{2\pi}|\widetilde{\mathbf E}(\omega)|^2.
\end{equation*}
实场满足 $\widetilde{\mathbf E}(-\omega)=\widetilde{\mathbf E}^{*}(\omega)$，因此正频部分为
\begin{equation*}
 \int\dd t\,|\mathbf E(t)|^2
 =\frac1\pi\int_0^\infty\dd\omega\,
 |\widetilde{\mathbf E}(\omega)|^2.
\end{equation*}
正文的远场源--场关系是
\begin{equation*}
 \widetilde{\mathbf E}_{\rm far}
 =\ii\mu_0\omega\frac{\ee^{\ii\kappa_0r}}{r}
 \mathbf A_\infty,
 \qquad \kappa_0=\omega/c.
\end{equation*}
代入 Poynting 谱后，$r^{-2}$ 与球面面积因子相消，得到
\begin{equation*}
 \frac{\dd^2 W_{\rm rad}}{\dd\omega\dd\Omega}
 =\frac{\mu_0\omega^2}{\pi c}|\mathbf A_\infty|^2,
\end{equation*}
即正文\eqnrefs{eq:far_field_radiated_energy_green_si}。正文\eqnrefs{eq:cl_photon_probability_far_green} 随后只是把这一能量谱除以单光子能量 $\hbar\omega$ 的定义转换，不需要另行引入动力学近似。
\end{derivation}

\begin{derivation}[自由空间电流的光子辐射概率]
真空标量 Helmholtz Green 函数
$g_0(\mathbf r-\mathbf r')=\ee^{\ii\kappa_0|\mathbf r-\mathbf r'|}/[4\pi|\mathbf r-\mathbf r'|]$ 的远场系数由严格极限定义：
\begin{equation*}
\lim_{r\to\infty}
4\pi r\,\ee^{-\ii\kappa_0r}
 g_0(r\hat{\mathbf n}-\mathbf r')
=\ee^{-\ii\kappa_0\hat{\mathbf n}\cdot\mathbf r'}.
\end{equation*}
电磁 Green 张量在同一极限只保留横向投影，因此
\[
\mathbf F_{\infty,0}
(\hat{\mathbf n},\mathbf r';\omega)
=\frac1{4\pi}
(\mathbf I-\hat{\mathbf n}\hat{\mathbf n})
\ee^{-\ii\kappa_0\hat{\mathbf n}\cdot\mathbf r'}.
\]
代入 $\mathbf A_\infty$ 定义：
\begin{align*}
\mathbf A_{\infty,0}
&=\int\dd^3 r'\,
\mathbf F_{\infty,0}\cdot\widetilde{\mathbf J}(\mathbf r',\omega)\\
&=\frac1{4\pi}
(\mathbf I-\hat{\mathbf n}\hat{\mathbf n})
\int\dd^3 r'\,\widetilde{\mathbf J}(\mathbf r',\omega)
\ee^{-\ii\kappa_0\hat{\mathbf n}\cdot\mathbf r'}\\
&=\frac1{4\pi}
(\mathbf I-\hat{\mathbf n}\hat{\mathbf n})
\widetilde{\mathbf J}(\kappa_0\hat{\mathbf n},\omega).
\end{align*}
因此
\[
|\mathbf A_{\infty,0}|^2
=\frac1{16\pi^2}
\left|(\mathbf I-\hat{\mathbf n}\hat{\mathbf n})
\widetilde{\mathbf J}\right|^2.
\]
代入\eqnrefs{eq:cl_photon_probability_far_green} 后，$1/\pi$ 与 $1/(16\pi^2)$ 合并成 $1/(16\pi^3)$，得到\eqnrefs{eq:free_space_current_radiation_probability}。
\end{derivation}

\begin{derivation}[能量损失、吸收与远场辐射的平衡]

先由外部电流对诱导场所做的负功求总损失。时域表达式为
\begin{equation*}
 W_{\rm loss}
 =-\int\dd t\dd^3 r\,
 \mathbf J(\mathbf r,t)\cdot\mathbf E(\mathbf r,t).
\end{equation*}
利用实场的正负频共轭关系，正频谱写成
\begin{equation*}
 \frac{\dd W_{\rm loss}}{\dd\omega}
 =-\frac1\pi\operatorname{Re}
 \int\dd^3 r\,
 \widetilde{\mathbf J}^{*}(\mathbf r,\omega)
 \cdot\widetilde{\mathbf E}(\mathbf r,\omega).
\end{equation*}
代入正文\eqnrefs{eq:eels_source_field_green_dp68}，
\begin{align*}
 \frac{\dd W_{\rm loss}}{\dd\omega}
 &=-\frac1\pi\operatorname{Re}
 \left[\ii\mu_0\omega
 \langle\widetilde J|G^R|\widetilde J\rangle\right]\\
 &=\frac{\mu_0\omega}{\pi}
 \langle\widetilde J|\operatorname{Im}G^R|\widetilde J\rangle,
\end{align*}
即正文\eqnrefs{eq:eels_energy_loss_green_si}。

再计算局域介质中的不可逆吸收。正文\eqnrefs{eq:material_polarization_current_dp68} 给出
\begin{equation*}
 \widetilde{\mathbf J}_{\rm mat}
 =-\ii\omega\varepsilon_0
 (\boldsymbol\varepsilon_r-\mathbf I)\widetilde{\mathbf E}.
\end{equation*}
取正频周期平均功，
\begin{align*}
 \frac{\dd W_{\rm abs}}{\dd\omega}
 &=\frac1\pi\int_{V_{\rm mat}}\dd^3 r\,
 \operatorname{Re}
 \left[\widetilde{\mathbf J}_{\rm mat}\cdot
 \widetilde{\mathbf E}^{*}\right]\\
 &=\frac{\omega\varepsilon_0}{\pi}
 \int_{V_{\rm mat}}\dd^3 r\,
 \widetilde{\mathbf E}^{*}\cdot
 \operatorname{Im}\boldsymbol\varepsilon_r
 \cdot\widetilde{\mathbf E},
\end{align*}
得到正文\eqnrefs{eq:material_absorbed_energy_spectrum_si}。

最后证明总损失如何严格分成远场辐射与材料吸收。取一个包围全部耗散材料、外边界位于无损真空的控制体 $V$。按全文 $\ee^{-\ii\omega t}$ 约定，
\begin{equation*}
 \nabla\times\mathbf E=\ii\omega\mathbf B,
 \qquad
 \nabla\times\mathbf H=\mathbf J_{\rm ext}-\ii\omega\mathbf D,
\end{equation*}
所以共轭方程为
\begin{equation*}
 \nabla\times\mathbf H^*=\mathbf J_{\rm ext}^*+\ii\omega\mathbf D^*.
\end{equation*}
代入恒等式
\begin{equation*}
 \nabla\cdot(\mathbf E\times\mathbf H^*)
 =\mathbf H^*\cdot(\nabla\times\mathbf E)
 -\mathbf E\cdot(\nabla\times\mathbf H^*)
\end{equation*}
得到复 Poynting 方程
\begin{equation*}
 \nabla\cdot(\mathbf E\times\mathbf H^*)
 =\ii\omega\mathbf H^*\cdot\mathbf B
 -\mathbf E\cdot\mathbf J_{\rm ext}^*
 -\ii\omega\mathbf E\cdot\mathbf D^*.
\end{equation*}
介质只含电损耗并取 $\mathbf B=\mu_0\mathbf H$，第一项为纯虚数，对实部没有贡献。再写
\begin{equation*}
 \mathbf D=\varepsilon_0\boldsymbol\varepsilon_r\mathbf E,
 \qquad
 X\equiv\mathbf E^*\cdot\boldsymbol\varepsilon_r\mathbf E,
\end{equation*}
则 $\mathbf E\cdot\mathbf D^*=\varepsilon_0X^*$，从而
\begin{equation*}
 \operatorname{Re}[-\ii\omega\mathbf E\cdot\mathbf D^*]
 =-\omega\varepsilon_0\operatorname{Im}X
 =-\omega\varepsilon_0
 \mathbf E^*\cdot\operatorname{Im}\boldsymbol\varepsilon_r\cdot\mathbf E.
\end{equation*}
因此
\begin{align*}
 \operatorname{Re}\nabla\cdot(\mathbf E\times\mathbf H^*)
 ={}&-\operatorname{Re}(\mathbf J_{\rm ext}^*\cdot\mathbf E)\\
 &-\omega\varepsilon_0
 \mathbf E^*\cdot\operatorname{Im}\boldsymbol\varepsilon_r\cdot\mathbf E.
\end{align*}
对 $V$ 积分并用 Gauss 定理，
\begin{align*}
 \operatorname{Re}\oint_{\partial V}\dd\mathbf S\cdot
 (\mathbf E\times\mathbf H^*)
 ={}&-\operatorname{Re}\int_V\dd^3r\,
 \mathbf J_{\rm ext}^*\cdot\mathbf E\\
 &-\omega\varepsilon_0\int_{V_{\rm mat}}\dd^3r\,
 \mathbf E^*\cdot\operatorname{Im}\boldsymbol\varepsilon_r\cdot\mathbf E.
\end{align*}
全文的正频 Fourier 归一化给每一项共同的 $1/\pi$。于是右边第一项正是前面得到的 $\dd W_{\rm loss}/\dd\omega$，第二项是 $-\dd W_{\rm abs}/\dd\omega$；把 $\partial V$ 推到无损真空远场后，左边就是 $\dd W_{\rm rad}/\dd\omega$。因此逐频有
\begin{equation*}
 \frac{\dd W_{\rm loss}}{\dd\omega}
 =\frac{\dd W_{\rm rad}}{\dd\omega}
 +\frac{\dd W_{\rm abs}}{\dd\omega},
\end{equation*}
即正文\eqnrefs{eq:eels_radiative_absorptive_energy_balance}。这里的分解来自同一个 Maxwell 能量守恒恒等式，而不是把 EELS 与 CL 当成两套独立动力学。
\end{derivation}

\begin{derivation}[平面反射 Green 张量的纵向分量]
真空侧平面波的单位 TE 偏振定义为

\[
\hat{\mathbf s}
=\frac{\hat{\mathbf z}\times\mathbf \kappa_\parallel}{\kappa_\parallel}.
\]
若
\[
\mathbf \kappa_\parallel=\kappa_x\hat{\mathbf x}+\kappa_y\hat{\mathbf y},
\]
则
\[
\hat{\mathbf s}
=\frac{-\kappa_y\hat{\mathbf x}+\kappa_x\hat{\mathbf y}}{\kappa_\parallel},
\]
因此
\[
|\hat{\mathbf x}\cdot\hat{\mathbf s}|^2
=\frac{\kappa_y^2}{\kappa_\parallel^2}.
\]

真空 TM 偏振取
\[
\hat{\mathbf p}_1^\pm
=\frac{\kappa_\parallel\hat{\mathbf z}
\mp \kappa_{z1}\hat{\boldsymbol\kappa}_\parallel}{\kappa_0},
\qquad
\hat{\boldsymbol\kappa}_\parallel
=\frac{\mathbf \kappa_\parallel}{\kappa_\parallel}.
\]
所以
\begin{align*}
\hat{\mathbf x}\cdot\hat{\mathbf p}_1^+
&=-\frac{\kappa_{z1}}{\kappa_0}\frac{\kappa_x}{\kappa_\parallel},\\
\hat{\mathbf p}_1^-\cdot\hat{\mathbf x}
&=+\frac{\kappa_{z1}}{\kappa_0}\frac{\kappa_x}{\kappa_\parallel},
\end{align*}
从而
\[
(\hat{\mathbf x}\cdot\hat{\mathbf p}_1^+)
(\hat{\mathbf p}_1^-\cdot\hat{\mathbf x})
=-\frac{\kappa_{z1}^2}{\kappa_0^2}
\frac{\kappa_x^2}{\kappa_\parallel^2}.
\]
亚光速同步条件给 $\kappa_\parallel>\kappa_0$，故
\[
\kappa_{z1}=\ii\alpha_1,
\qquad
-\kappa_{z1}^2=\alpha_1^2.
\]
于是 TM 投影化为
\[
\frac{\alpha_1^2}{\kappa_0^2}
\frac{\kappa_x^2}{\kappa_\parallel^2}.
\]
Weyl 散射 Green 张量的公共传播因子为
\[
\frac{\ii}{2\kappa_{z1}}
\ee^{\ii \kappa_{z1}(z+z')}.
\]
取源点和观察点都在电子轨迹 $z=z'=b$：
\begin{align*}
\frac{\ii}{2\kappa_{z1}}\ee^{2\ii \kappa_{z1}b}
&=\frac{\ii}{2\ii\alpha_1}\ee^{-2\alpha_1 b}\\
&=\frac{\ee^{-2\alpha_1 b}}{2\alpha_1}.
\end{align*}
TE 投影乘该因子给
\[
\ee^{-2\alpha_1 b}\frac{\kappa_y^2}{2\alpha_1 \kappa_\parallel^2}r_s,
\]
TM 投影乘该因子给
\[
\ee^{-2\alpha_1 b}
\frac{\alpha_1 \kappa_x^2}{2\kappa_0^2\kappa_\parallel^2}r_p.
\]
两项相加得到\eqnrefs{eq:planar_aloof_gxx_explicit}。
\end{derivation}

\begin{derivation}[平行掠射电子的 Fresnel 损失积分]
三维 Green 张量满足

\[
[\nabla\times\nabla\times-\kappa_0^2\varepsilon]\mathbf G^R
=\mathbf I\delta^{(3)}(\mathbf r-\mathbf r'),
\]
因此 $[\mathbf G^R]={\rm m^{-1}}$。Weyl 表示定义为
\[
\mathbf G^R
=\int\frac{\dd^2 \kappa_\parallel}{(2\pi)^2}
\ee^{\ii\mathbf \kappa_\parallel\cdot(\mathbf R-\mathbf R')}
\mathbf G_{{\rm sc},W}^R.
\]
因为 $[\dd^2 \kappa_\parallel]={\rm m^{-2}}$，必须有
\[
[\mathbf G_{{\rm sc},W}^R]={\rm m}.
\]
因此 $\dd \kappa_y\,G_{{\rm sc},W;xx}^R$ 无量纲。

前因子单位为
\begin{align*}
\left[\frac{\mu_0e^2}{\hbar}\right]
&=\frac{({\rm kg\,m\,C^{-2}})({\rm C^2})}
{{\rm J\,s}}\\
&=\frac{{\rm kg\,m}}{{\rm kg\,m^2\,s^{-1}}}\\
&={\rm s\,m^{-1}}.
\end{align*}
所以
\[
\left[\frac1L\frac{\dd P}{\dd\omega}\right]
={\rm s\,m^{-1}}.
\]
若换变量 $E_{\rm loss}=\hbar\omega$，
\[
\frac{\dd\omega}{\dd E_{\rm loss}}=\frac1\hbar,
\]
因此
\[
\frac1L\frac{\dd P}{\dd E_{\rm loss}}
=\frac1\hbar\frac1L\frac{\dd P}{\dd\omega}.
\]
\end{derivation}

\begin{derivation}[Drude 表面模与电子相位同步]
无损 Drude 介电函数为

\[
\varepsilon_2(\omega)=1-\frac{\omega_p^2}{\omega^2}.
\]
真空--介质界面 TM 表面模色散满足
\[
\frac{\kappa_{\rm sp}}{\kappa_0}
=\sqrt{\frac{\varepsilon_2}{1+\varepsilon_2}}.
\]
在 $\kappa_y=0$ 时，电子源的同步条件是
\[
\kappa_{\rm sp}=\frac{\omega}{v}
=\frac{\kappa_0}{\beta_{\rm rel}}.
\]
平方并消去 $\kappa_0$：
\[
\frac1{\beta_{\rm rel}^2}
=\frac{\varepsilon_2}{1+\varepsilon_2}.
\]
交叉相乘：
\begin{align*}
1+\varepsilon_2
&=\beta_{\rm rel}^2\varepsilon_2,\\
1
&=(\beta_{\rm rel}^2-1)\varepsilon_2,\\
\varepsilon_2
&=-\frac1{1-\beta_{\rm rel}^2}.
\end{align*}
再代入 Drude 形式：
\begin{align*}
1-\frac{\omega_p^2}{\omega^2}
&=-\frac1{1-\beta_{\rm rel}^2},\\
\frac{\omega_p^2}{\omega^2}
&=1+\frac1{1-\beta_{\rm rel}^2}\\
&=\frac{2-\beta_{\rm rel}^2}{1-\beta_{\rm rel}^2}.
\end{align*}
所以
\[
\omega^2
=\omega_p^2\frac{1-\beta_{\rm rel}^2}{2-\beta_{\rm rel}^2}.
\]
取正频率根得到\eqnrefs{eq:drude_surface_pole_electron_sync_frequency}。
\end{derivation}

